\documentclass[a4paper,12pt]{article}

% Пакеты для русского языка и кодировки
\usepackage[utf8]{inputenc}
\usepackage[T2A]{fontenc}
\usepackage[russian]{babel}

% Пакеты для математики
\usepackage{amsmath}
\usepackage{amssymb}
\usepackage{amsthm}

% Пакеты для оформления
\usepackage{geometry}
\geometry{left=2.5cm, right=2.5cm, top=2.5cm, bottom=2.5cm}
\usepackage{hyperref}
\usepackage{enumitem}
\usepackage{booktabs}
\usepackage{array}

% Настройка заголовков
\title{2.3.5 Материалы для изучения по научной специальности \\
Математическое и программное обеспечение вычислительных систем, комплексов и компьютерных сетей \\
Гусев Виталий Евгеньевич 2.3.5 2026}
\date{}

\begin{document}

\maketitle

\tableofcontents
\newpage

\setcounter{section}{2} % чтобы раздел имел номер 3
\section{Раздел 3. Вычислительные машины, системы и сети. Языки программирования.}

\noindent\textit{Основные источники: Степанов А.Н. Архитектура вычислительных систем и компьютерных сетей. СПб.: Питер, 2007. Часть III «Введение в архитектуру компьютерных сетей», гл. 19–26, с. 432–482.}

\subsection{Назначение, архитектура и принципы построения информационно-вычислительных сетей (ИВС)}

\noindent\textit{Основные источники: Степанов А.Н. Архитектура вычислительных систем и компьютерных сетей. СПб.: Питер, 2007. Часть III «Введение в архитектуру компьютерных сетей», гл. 19–26, с. 432–482.}

\subsubsection{Понятие и назначение ИВС}

\paragraph{Определение информационно-вычислительной сети.}
Информационно-вычислительная сеть (ИВС) представляет собой совокупность компьютеров, терминалов, периферийных устройств, средств передачи данных и сетевого программного обеспечения, объединённых в единую систему для совместного использования информационных, аппаратных и программных ресурсов, а также для распределённой обработки информации. В рамках учебника Степанова А.Н. ИВС рассматривается как компьютерная сеть, ориентированная на автоматизацию процессов приёма, хранения, обработки и передачи информации. Отдельный компьютер при этом остаётся универсальным средством обработки данных, а сеть обеспечивает качественно новый уровень — коллективную и распределённую работу множества устройств.
\textit{(Источник: Степанов, ч. III, гл. 19–20, с. 432–450.)}

\paragraph{Основное назначение ИВС.}
Назначение ИВС состоит в обеспечении обмена данными между удалёнными устройствами и организации коллективного доступа к ресурсам. К типовым задачам, решаемым с помощью ИВС, относятся:
\begin{itemize}
\item совместное использование аппаратных ресурсов — процессоров, накопителей, принтеров, сканеров, специализированного измерительного и исполнительного оборудования;
\item совместное использование программных и информационных ресурсов — файловых хранилищ, баз данных, прикладных программ, электронных библиотек, web-сервисов;
\item распределённая обработка данных, при которой вычислительная задача разбивается на части и решается несколькими компьютерами сети;
\item оперативный обмен сообщениями и документами — электронная почта, файлообмен, видеоконференции, мессенджеры;
\item централизованное хранение и резервирование данных, повышение надёжности и доступности информационных сервисов;
\item удалённое управление оборудованием, мониторинг и администрирование распределённых систем;
\item организация доступа к глобальным информационным ресурсам, включая Интернет.
\end{itemize}
Таким образом, ИВС обеспечивает переход от автономной обработки данных к распределённой и коллективной работе, что существенно повышает эффективность использования вычислительной техники.
\textit{(Источник: Степанов, ч. III, гл. 19–20, с. 432–450.)}

\paragraph{Основные эффекты от построения ИВС.}
Построение ИВС даёт следующие основные эффекты:
\begin{itemize}
\item \emph{ресурсный} — дорогостоящее оборудование и программное обеспечение используются коллективно, что снижает суммарные затраты;
\item \emph{вычислительный} — появляется возможность распараллеливания и распределения нагрузки между несколькими компьютерами;
\item \emph{информационный} — обеспечивается быстрый доступ к распределённым данным и их совместное использование;
\item \emph{организационный} — упрощается централизованное управление, мониторинг и администрирование;
\item \emph{надёжностный} — резервирование каналов, узлов и данных повышает устойчивость системы к отказам;
\item \emph{коммуникационный} — обеспечивается оперативное взаимодействие пользователей и устройств.
\end{itemize}
\textit{(Источник: Степанов, ч. III, гл. 19–26, с. 432–482.)}

\subsubsection{Архитектура ИВС}

\paragraph{Понятие архитектуры сети.}
Под архитектурой ИВС понимается совокупность принципов построения и взаимодействия аппаратных и программных средств сети, включающая физическую среду передачи данных, состав и топологию узлов, методы доступа к среде, методы коммутации, систему адресации, набор протоколов и структуру программного обеспечения. Архитектура сети определяет её функциональные возможности, масштабируемость, надёжность, стоимость и область применения.
\textit{(Источник: Степанов, ч. III, гл. 19–26, с. 432–482.)}

\paragraph{Компоненты архитектуры ИВС.}
В архитектуре ИВС принято выделять несколько взаимосвязанных компонентов, представленных в табл. \ref{tab:ivs-components}.

\begin{table}[h]
\centering
\caption{Основные компоненты архитектуры ИВС}
\label{tab:ivs-components}
\begin{tabular}{p{3.2cm}p{10.2cm}}
\toprule
\textbf{Компонент} & \textbf{Назначение и содержание} \\
\midrule
Линии связи & Физическая среда передачи сигналов: витая пара, коаксиальный кабель, оптоволокно, радиоканалы, спутниковые и инфракрасные каналы \\
Узлы сети & Компьютеры, серверы, рабочие станции, терминалы, сетевые принтеры и другие устройства, подключённые к сети \\
Сетевое оборудование & Сетевые адаптеры, повторители, концентраторы, мосты, коммутаторы, маршрутизаторы, шлюзы, модемы, точки доступа \\
Сетевое ПО & Сетевые операционные системы, драйверы, протоколы, службы, прикладные программы обмена и доступа к ресурсам \\
Адресация и маршрутизация & Правила назначения адресов узлам и определение путей передачи данных между ними \\
Протоколы & Согласованные правила представления, передачи, обработки и интерпретации данных \\
Топология & Физическая и логическая структура связей между узлами сети \\
\bottomrule
\end{tabular}
\end{table}

\textit{(Источник: Степанов, ч. III, гл. 19–22, с. 432–461.)}

\paragraph{Линии связи и организация передачи данных.}
Линии связи образуют физическую основу сети. Они характеризуются типом среды, пропускной способностью, помехоустойчивостью, задержкой распространения сигнала, затуханием и стоимостью. В гл. 19 учебника рассматриваются режимы передачи сообщений, параллельная и последовательная передача, способы представления кодов, а также методы обнаружения и исправления ошибок. Параллельная передача предполагает одновременную передачу нескольких битов по нескольким линиям, последовательная — поочерёдную передачу битов по одной линии. Для повышения достоверности применяются контрольные разряды, коды с обнаружением и исправлением ошибок, повторная передача искажённых блоков.
\textit{(Источник: Степанов, ч. III, гл. 19, с. 432–443.)}

\paragraph{Топология сети.}
Топология ИВС характеризует способ соединения узлов линиями связи. Различают физическую и логическую топологию. Физическая топология описывает реальное расположение кабелей и устройств, а логическая — структуру потоков данных и адресации. К основным топологиям относятся:
\begin{itemize}
\item \emph{шина} — все узлы подключены к общему каналу передачи; достоинства — простота и экономичность, недостатки — зависимость работы сети от состояния общего канала и ограниченная масштабируемость;
\item \emph{звезда} — каждый узел подключён к центральному устройству (концентратору или коммутатору); достоинства — удобство управления и локализация неисправностей, недостатки — зависимость от центрального устройства;
\item \emph{кольцо} — узлы образуют замкнутый контур, данные передаются по кольцу; используются маркерные методы доступа;
\item \emph{дерево} — иерархическое соединение сегментов и узлов; характерно для структурированных кабельных систем;
\item \emph{полносвязная и смешанная топологии} — обеспечивают высокую надёжность и альтернативные пути, но требуют большего количества соединений и стоимости.
\end{itemize}
Выбор топологии определяется требованиями к производительности, надёжности, стоимости и условиям прокладки линий связи.
\textit{(Источник: Степанов, ч. III, гл. 20, с. 444–450.)}

\paragraph{Методы доступа к среде передачи.}
В ИВС применяются различные методы доступа к разделяемой среде передачи данных. Основными являются:
\begin{itemize}
\item \emph{случайный доступ} — узел начинает передачу при необходимости, а конфликты разрешаются специальными процедурами; типичный пример — метод CSMA/CD, использовавшийся в классическом Ethernet;
\item \emph{маркерный доступ} — право на передачу передаётся по сети специальным маркером; такой подход применяется в Token Ring и FDDI и обеспечивает детерминированное время доступа.
\end{itemize}
Выбор метода доступа влияет на пропускную способность, задержки и устойчивость сети при высокой нагрузке.
\textit{(Источник: Степанов, ч. III, гл. 23, с. 462–466.)}

\paragraph{Методы коммутации.}
Для передачи данных между узлами сети используются методы коммутации:
\begin{itemize}
\item \emph{коммутация каналов} — на время сеанса связи выделяется физический или логический канал с постоянной пропускной способностью;
\item \emph{коммутация сообщений} — сообщение передаётся целиком с промежуточным хранением на узлах;
\item \emph{коммутация пакетов} — сообщение разбивается на пакеты, которые независимо передаются по сети и собираются в узле-получателе.
\end{itemize}
В современных ИВС доминирует коммутация пакетов, обеспечивающая эффективное использование каналов и гибкость маршрутизации.
\textit{(Источник: Степанов, ч. III, гл. 24, с. 467–471.)}

\paragraph{Адресация и маршрутизация.}
Для идентификации узлов и доставки данных в ИВС используются многоуровневые адреса. На канальном уровне применяются MAC-адреса сетевых адаптеров, на сетевом — IP-адреса, на прикладном — доменные имена. Маршрутизаторы анализируют адреса пакетов и выбирают путь передачи по таблицам маршрутизации; для обмена маршрутной информацией применяются протоколы маршрутизации. Иерархическая адресация и маршрутизация особенно важны в глобальных сетях, где требуется объединять множество разнородных сетей.
\textit{(Источник: Степанов, ч. III, гл. 20, 25–26, с. 444–450, 472–482.)}

\paragraph{Сетевое оборудование.}
К техническим средствам ИВС относятся сетевые адаптеры, повторители, концентраторы, мосты, коммутаторы, маршрутизаторы, шлюзы, модемы и точки доступа. Повторители и концентраторы работают на физическом уровне, мосты и коммутаторы — на канальном, маршрутизаторы — на сетевом, шлюзы — на прикладном и вышележащих уровнях. Коммутаторы обеспечивают передачу кадров между портами и разделение доменов коллизий; маршрутизаторы соединяют разные сети и разделяют широковещательные домены; шлюзы выполняют трансляцию протоколов при объединении разнородных сетей.
\textit{(Источник: Степанов, ч. III, гл. 21–22, с. 451–461.)}

\paragraph{Программное обеспечение ИВС.}
Программную основу ИВС образуют сетевые операционные системы, драйверы сетевых адаптеров, стеки протоколов, файловые и печатные службы, службы каталогов, DNS, DHCP, почтовые и web-серверы, средства удалённого доступа и администрирования. Наиболее распространён стек TCP/IP, включающий протоколы IP, TCP, UDP, ICMP, ARP и прикладные протоколы HTTP, FTP, SMTP, DNS, DHCP и др. Протоколы определяют форматы пакетов, правила адресации, установления соединения, маршрутизации и обработки ошибок.
\textit{(Источник: Степанов, ч. III, гл. 25–26, с. 472–482.)}

\paragraph{Многоуровневая модель OSI.}
Архитектура сетевого взаимодействия описывается многоуровневой моделью OSI, включающей семь уровней. Каждый уровень выполняет определённый набор функций и предоставляет сервисы вышележащему уровню, используя услуги нижележащего. Такая декомпозиция обеспечивает открытость, совместимость и модульность сетевых решений. Соответствие уровней и их функций приведено в табл. \ref{tab:osi-levels}.

\begin{table}[h]
\centering
\caption{Уровни модели OSI и их основные функции}
\label{tab:osi-levels}
\begin{tabular}{p{1.2cm}p{4.2cm}p{7.8cm}}
\toprule
\textbf{№} & \textbf{Уровень} & \textbf{Основные функции} \\
\midrule
7 & Прикладной & Доступ к сетевым службам, прикладные протоколы, интерфейс пользователя \\
6 & Представления данных & Кодирование, сжатие, шифрование, преобразование форматов \\
5 & Сеансовый & Установление, поддержание и завершение сеансов связи \\
4 & Транспортный & Надёжная доставка данных, сегментация, управление потоком и ошибками \\
3 & Сетевой & Адресация, маршрутизация, коммутация пакетов между сетями \\
2 & Канальный & Формирование кадров, MAC-адресация, доступ к среде, обнаружение ошибок \\
1 & Физический & Передача битов по среде, кодирование сигналов, характеристики линий связи \\
\bottomrule
\end{tabular}
\end{table}

\textit{(Источник: Степанов, ч. III, гл. 26, с. 475–482.)}

\subsubsection{Принципы построения ИВС}

\paragraph{Основные принципы.}
К основным принципам построения современных ИВС относятся:
\begin{itemize}
\item \emph{открытость и стандартизация} — использование международных стандартов на протоколы, интерфейсы и форматы данных, что обеспечивает совместимость оборудования разных производителей;
\item \emph{многоуровневость и декомпозиция} — разбиение сложной задачи сетевого взаимодействия на относительно независимые уровни и подсистемы;
\item \emph{модульность} — построение сети из унифицированных компонентов, что упрощает проектирование, модернизацию и расширение;
\item \emph{совместимость и интероперабельность} — способность различных сетей и устройств совместно работать и обмениваться данными;
\item \emph{масштабируемость} — возможность увеличения числа узлов и протяжённости сети без принципиальной перестройки архитектуры;
\item \emph{надёжность и отказоустойчивость} — резервирование каналов и оборудования, альтернативные маршруты, автоматическое восстановление после сбоев;
\item \emph{эффективность} — обеспечение требуемой пропускной способности, минимальных задержек и рационального использования ресурсов;
\item \emph{безопасность} — защита информации от несанкционированного доступа, аутентификация, шифрование, сегментация сети, межсетевые экраны;
\item \emph{управляемость} — наличие средств мониторинга, диагностики, конфигурирования и администрирования сети;
\item \emph{иерархичность адресации и маршрутизации} — использование многоуровневых адресов и маршрутных схем, характерных для локальных и глобальных сетей.
\end{itemize}
\textit{(Источник: Степанов, ч. III, гл. 19–26, с. 432–482.)}

\paragraph{Физическая и логическая структуризация.}
Важным принципом построения ИВС является структуризация. Физическая структуризация предполагает разделение сети на сегменты с помощью повторителей, концентраторов, коммутаторов и структурированной кабельной системы. Логическая структуризация связана с разделением сети на логические подсети, VLAN, зоны адресации и маршрутизации. Структуризация позволяет повысить производительность, локализовать трафик, упростить управление и повысить безопасность.
\textit{(Источник: Степанов, ч. III, гл. 22, с. 457–461.)}

\paragraph{Принцип пакетной передачи и коммутации.}
Современные ИВС строятся на основе пакетной передачи данных. Сообщение разбивается на пакеты, каждый из которых содержит служебную информацию — адреса отправителя и получателя, номер пакета, контрольную сумму и другие данные. Пакеты независимо передаются по сети и собираются в узле-получателе. Пакетная коммутация обеспечивает эффективное использование каналов, возможность альтернативной маршрутизации и устойчивость к перегрузкам. Принцип пакетной передачи тесно связан с принципами буферизации, маршрутизации и управления потоком.
\textit{(Источник: Степанов, ч. III, гл. 24, с. 467–471.)}

\paragraph{Принцип открытых систем.}
Принцип открытых систем означает, что архитектура и протоколы сети должны быть описаны в общедоступных стандартах, а взаимодействие компонентов должно обеспечиваться независимо от их производителя. Это позволяет объединять в одной сети оборудование и программное обеспечение различных фирм, заменять отдельные компоненты без перестройки всей сети и создавать переносимые прикладные программы. Модель OSI является теоретической основой построения открытых сетей.
\textit{(Источник: Степанов, ч. III, гл. 26, с. 475–482.)}

\paragraph{Принципы надёжности и безопасности.}
Надёжность ИВС обеспечивается резервированием каналов и оборудования, альтернативной маршрутизацией, контролем ошибок, буферизацией и повторной передачей. Безопасность достигается аутентификацией пользователей и устройств, разграничением доступа, шифрованием трафика, сегментацией сети, применением межсетевых экранов, систем обнаружения и предотвращения вторжений, а также организационными мерами. Эти принципы особенно важны при объединении локальных и глобальных сетей.
\textit{(Источник: Степанов, ч. III, гл. 21–26, с. 451–482.)}

\subsubsection{Краткое резюме подраздела}
ИВС предназначена для совместного использования ресурсов и распределённой обработки информации. Архитектура ИВС включает линии связи, узлы, сетевое оборудование, программное обеспечение, систему адресации, протоколы и топологию. Ключевыми принципами построения являются открытость, многоуровневость, модульность, совместимость, масштабируемость, надёжность, эффективность, безопасность, управляемость, иерархичность адресации и структуризация. Для описания взаимодействия используется модель OSI; передача данных организуется методами коммутации, из которых основным является коммутация пакетов. Совместное применение этих принципов обеспечивает построение масштабируемых, надёжных и совместимых информационно-вычислительных сетей.
\textit{(Обобщение по: Степанов, ч. III, гл. 19–26, с. 432–482.)}

\subsection{Локальные и глобальные ИВС, технические и программные средства объединения различных сетей}

\noindent\textit{Основные источники: Степанов А.Н. Архитектура вычислительных систем и компьютерных сетей. СПб.: Питер, 2007. Часть III «Введение в архитектуру компьютерных сетей», гл. 19–26, с. 432–482.}

\subsubsection{Локальные вычислительные сети}

\paragraph{Определение и особенности ЛВС.}
Локальная вычислительная сеть (ЛВС, LAN, Local Area Network) — это сеть, охватывающая ограниченную территорию: помещение, здание, комплекс зданий, предприятие или учреждение. Для ЛВС характерны:
\begin{itemize}
\item небольшие расстояния между узлами;
\item высокая скорость передачи данных;
\item низкий уровень ошибок;
\item принадлежность одной организации;
\item относительно простая топология и однородное оборудование;
\item использование разделяемых ресурсов в пределах локального коллектива.
\end{itemize}
ЛВС обеспечивает совместное использование файлов, принтеров, баз данных, прикладных программ и доступ к локальным серверам. Именно ЛВС исторически стали первой формой объединения компьютеров, позволившей перейти от автономной обработки данных к коллективной работе.
\textit{(Источник: Степанов, ч. III, гл. 20, с. 444–450.)}

\paragraph{Топологии ЛВС.}
В ЛВС применяются топологии «шина», «звезда», «кольцо», «дерево» и их комбинации. Краткая характеристика основных топологий приведена в табл. \ref{tab:lan-topologies}.

\begin{table}[h]
\centering
\caption{Основные топологии ЛВС}
\label{tab:lan-topologies}
\begin{tabular}{p{2.2cm}p{5.5cm}p{5.5cm}}
\toprule
\textbf{Топология} & \textbf{Достоинства} & \textbf{Недостатки} \\
\midrule
Шина & Простота, экономичность, малое количество кабеля & Зависимость от общего канала, ограниченная масштабируемость, сложность локализации неисправностей \\
Звезда & Удобство управления, локализация неисправностей, простота подключения новых узлов & Зависимость от центрального устройства, дополнительные затраты на кабель \\
Кольцо & Детерминированное время доступа, отсутствие коллизий при маркерном методе & Выход из строя одного узла может нарушить работу кольца, сложность реконфигурации \\
Дерево & Иерархичность, удобство структурирования, масштабируемость & Большая стоимость кабельной системы, зависимость от вышестоящих уровней \\
\bottomrule
\end{tabular}
\end{table}

Выбор топологии определяется требованиями к производительности, надёжности, стоимости и условиям прокладки линий связи.
\textit{(Источник: Степанов, ч. III, гл. 20, с. 444–450.)}

\paragraph{Технологии и среда передачи ЛВС.}
Классическими технологиями ЛВС являются Ethernet, Token Ring, FDDI; в современных сетях широко используются Fast Ethernet, Gigabit Ethernet, беспроводные сети Wi-Fi. В качестве среды передачи применяются витая пара, коаксиальный кабель и оптоволокно. Витая пара является наиболее распространённой средой для подключения рабочих станций; оптоволокно применяется на магистральных участках и для связи между зданиями, где требуются высокая пропускная способность и помехоустойчивость. Беспроводные технологии обеспечивают мобильность пользователей и упрощают развёртывание сети в труднодоступных местах.
\textit{(Источник: Степанов, ч. III, гл. 19–20, 25, с. 432–450, 472–474.)}

\paragraph{Оборудование ЛВС.}
К типовому оборудованию ЛВС относятся сетевые адаптеры, повторители, концентраторы, мосты, коммутаторы, точки беспроводного доступа и маршрутизаторы для связи с другими сетями. Сетевой адаптер обеспечивает подключение компьютера к среде передачи и формирование кадров. Концентратор объединяет узлы в общий домен коллизий и работает как многопортовый повторитель. Коммутатор обеспечивает передачу кадров между портами на канальном уровне и разделение доменов коллизий; он хранит таблицу MAC-адресов и направляет кадр только в тот порт, к которому подключён получатель. Маршрутизатор используется для объединения ЛВС с глобальными сетями и другими локальными сетями.
\textit{(Источник: Степанов, ч. III, гл. 21, с. 451–456.)}

\paragraph{Методы доступа в ЛВС.}
В ЛВС применяются случайный и маркерный методы доступа. Метод CSMA/CD, использовавшийся в классическом Ethernet, предполагает прослушивание среды перед передачей и обнаружение коллизий; при их возникновении передача повторяется через случайный интервал. Маркерный метод, применявшийся в Token Ring и FDDI, обеспечивает детерминированное время доступа за счёт передачи специального маркера. Современные коммутируемые ЛВС в значительной степени устраняют коллизии за счёт микросегментации и дуплексного режима.
\textit{(Источник: Степанов, ч. III, гл. 23, с. 462–466.)}

\subsubsection{Глобальные вычислительные сети}

\paragraph{Определение и особенности ГВС.}
Глобальная вычислительная сеть (ГВС, WAN, Wide Area Network) охватывает большие территории — регионы, страны, континенты. Для ГВС характерны:
\begin{itemize}
\item значительная протяжённость линий связи;
\item использование арендованных каналов, магистральных сетей и коммуникационного оборудования операторов связи;
\item гетерогенность оборудования и программного обеспечения;
\item сложная маршрутизация и глобальная адресация;
\item более высокие задержки и стоимость передачи по сравнению с ЛВС;
\item необходимость обеспечения надёжности и альтернативных маршрутов.
\end{itemize}
К ГВС относятся сети X.25, Frame Relay, ATM, SDH/SONET, MPLS, а также Интернет. ГВС обеспечивают связь между локальными сетями и удалёнными пользователями.
\textit{(Источник: Степанов, ч. III, гл. 20, 25, с. 444–450, 472–474.)}

\paragraph{Технологии глобальных сетей.}
Исторически первыми глобальными сетями были сети с коммутацией каналов и сообщений. В дальнейшем широкое распространение получили сети с коммутацией пакетов. Технология X.25 обеспечивала надёжную передачу по ненадёжным линиям за счёт развитых процедур контроля и исправления ошибок на каждом узле. Frame Relay упростила обработку кадров и уменьшила задержки, отказавшись от части процедур контроля. ATM использует ячейки фиксированной длины и обеспечивает высокое качество обслуживания. Технологии SDH/SONET образуют физическую магистральную основу, а MPLS обеспечивает эффективную коммутацию по меткам и поддержку виртуальных частных сетей. В настоящее время основой глобальных сетей является стек TCP/IP.
\textit{(Источник: Степанов, ч. III, гл. 25, с. 472–474.)}

\paragraph{Маршрутизация и адресация в ГВС.}
В глобальных сетях используется иерархическая адресация и маршрутизация. Маршрутизаторы анализируют адреса пакетов и выбирают путь передачи по таблицам маршрутизации; для обмена маршрутной информацией применяются протоколы маршрутизации. В Интернете используются IP-адресация, доменная система имён DNS, протоколы TCP/IP и система автономных систем. Иерархическая адресация позволяет агрегировать маршруты и сокращать размер таблиц маршрутизации, что особенно важно для магистральных сетей.
\textit{(Источник: Степанов, ч. III, гл. 20, 25–26, с. 444–450, 472–482.)}

\subsubsection{Сравнение ЛВС и ГВС}

Сравнительная характеристика ЛВС и ГВС приведена в табл. \ref{tab:lan-wan}.

\begin{table}[h]
\centering
\caption{Сравнительная характеристика ЛВС и ГВС}
\label{tab:lan-wan}
\begin{tabular}{p{3cm}p{4.5cm}p{4.5cm}}
\toprule
\textbf{Параметр} & \textbf{ЛВС} & \textbf{ГВС} \\
\midrule
Территория & Помещение, здание, предприятие & Регион, страна, континент \\
Собственность & Обычно одна организация & Операторы связи, множество организаций \\
Скорость передачи & Высокая & Ниже, зависит от канала \\
Задержки & Малые & Более высокие \\
Ошибки & Низкий уровень & Выше из-за протяжённости \\
Топология & Шина, звезда, кольцо, дерево & Смешанная, ячеистая, иерархическая \\
Оборудование & Коммутаторы, концентраторы, адаптеры & Маршрутизаторы, магистральные коммутаторы, модемы \\
Стоимость & Относительно низкая & Высокая \\
Стандарты & Ethernet, Wi-Fi, Token Ring & X.25, Frame Relay, ATM, MPLS, TCP/IP \\
\bottomrule
\end{tabular}
\end{table}

\textit{(Источник: Степанов, ч. III, гл. 19–26, с. 432–482.)}

\subsubsection{Технические средства объединения различных сетей}

\paragraph{Устройства объединения сетей.}
Объединение различных сетей обеспечивается сетевым оборудованием, работающим на разных уровнях модели OSI. Основные технические средства и их функции приведены в табл. \ref{tab:devices}.

\begin{table}[h]
\centering
\caption{Сетевое оборудование и уровень модели OSI}
\label{tab:devices}
\begin{tabular}{p{3.2cm}p{2.8cm}p{7.5cm}}
\toprule
\textbf{Устройство} & \textbf{Уровень OSI} & \textbf{Основная функция} \\
\midrule
Повторитель, концентратор & Физический & Регенерация сигнала, объединение сегментов, расширение среды передачи \\
Мост, коммутатор & Канальный & Фильтрация и передача кадров по MAC-адресам, разделение доменов коллизий \\
Маршрутизатор & Сетевой & Выбор маршрута, объединение сетей, IP-адресация, разделение широковещательных доменов \\
Шлюз & Прикладной и выше & Трансляция протоколов, объединение разнородных сетей \\
Модем, медиаконвертер & Физический & Согласование сигналов и сред передачи, подключение к линиям связи \\
Точка доступа & Канальный & Подключение беспроводных клиентов к проводной сети \\
\bottomrule
\end{tabular}
\end{table}

Повторитель (repeater) усиливает и регенерирует сигнал, позволяя увеличить протяжённость сегмента, но не фильтрует трафик. Концентратор (hub) представляет собой многопортовый повторитель и объединяет узлы в общий домен коллизий. Мост (bridge) работает на канальном уровне, соединяет сегменты и фильтрует кадры по MAC-адресам, разделяя домены коллизий. Коммутатор (switch) является многопортовым мостом, обеспечивает передачу кадров между портами, поддерживает таблицу MAC-адресов, VLAN, а часто выполняет и функции маршрутизации на третьем уровне. Маршрутизатор (router) работает на сетевом уровне, соединяет разные сети, выбирает маршрут по IP-адресам, разделяет широковещательные домены, поддерживает NAT и фильтрацию. Шлюз (gateway) выполняет трансляцию протоколов на прикладном и вышележащих уровнях, объединяя разнородные сети. Модемы, медиаконвертеры и точки доступа обеспечивают подключение к линиям связи и беспроводным средам.
\textit{(Источник: Степанов, ч. III, гл. 21–22, с. 451–461.)}

\paragraph{Линии связи и кабельная инфраструктура.}
Объединение сетей невозможно без линий связи. В современных ИВС используются витая пара, коаксиальный кабель, оптоволокно, радиоканалы, спутниковые каналы. Для построения структурированных кабельных систем применяются стандартизированные кабели, разъёмы, патч-панели и коммутационные шкафы. Оптоволокно обеспечивает высокую пропускную способность и помехоустойчивость на магистральных участках. Кабельная инфраструктура должна соответствовать стандартам, обеспечивать электромагнитную совместимость и возможность модернизации.
\textit{(Источник: Степанов, ч. III, гл. 19, 21–22, с. 432–443, 451–461.)}

\paragraph{Оборудование для связи ЛВС с ГВС.}
Для подключения локальных сетей к глобальным используются маршрутизаторы с интерфейсами WAN, модемы для выделенных и коммутируемых линий, устройства CSU/DSU, медиаконвертеры, а также оборудование операторов связи. На границе между локальной и глобальной сетью часто устанавливаются межсетевые экраны и прокси-серверы, обеспечивающие безопасность и контроль доступа.
\textit{(Источник: Степанов, ч. III, гл. 21–22, 25–26, с. 451–482.)}

\subsubsection{Программные средства объединения различных сетей}

\paragraph{Протоколы и стеки протоколов.}
Программной основой объединения сетей являются протоколы передачи данных. Наиболее распространён стек TCP/IP, включающий протоколы IP, TCP, UDP, ICMP, ARP и прикладные протоколы HTTP, FTP, SMTP, DNS, DHCP и др. Исторически применялись также IPX/SPX, NetBEUI и другие стеки. Протоколы определяют форматы пакетов, правила адресации, маршрутизации, установления соединения и обработки ошибок. Именно благодаря единому стеку протоколов становится возможным взаимодействие разнородных сетей и оборудования разных производителей.
\textit{(Источник: Степанов, ч. III, гл. 25–26, с. 472–482.)}

\paragraph{Сетевые операционные системы и службы.}
К программным средствам объединения сетей относятся сетевые операционные системы (Windows Server, Linux/Unix, Novell NetWare и др.), драйверы сетевых адаптеров, файловые и печатные службы, службы каталогов, DNS, DHCP, почтовые серверы, web-серверы, средства удалённого доступа. Они обеспечивают совместное использование ресурсов и управление сетью. Сетевые ОС реализуют протоколы, поддерживают сетевые службы и предоставляют администратору средства конфигурирования, мониторинга и разграничения доступа.
\textit{(Источник: Степанов, ч. III, гл. 21–26, с. 451–482.)}

\paragraph{Маршрутизация, туннелирование и трансляция.}
Для объединения различных сетей применяются:
\begin{itemize}
\item \emph{маршрутизация} — выбор пути передачи пакетов между сетями на основе таблиц маршрутизации и протоколов RIP, OSPF, BGP;
\item \emph{инкапсуляция} — вложение пакетов одной сети в пакеты другой сети при передаче через неоднородные среды;
\item \emph{туннелирование} — организация логического канала через промежуточную сеть, например VPN;
\item \emph{трансляция сетевых адресов NAT} — преобразование локальных адресов во внешние при выходе в глобальную сеть;
\item \emph{шлюзы и прокси-серверы} — согласование разнородных протоколов и прикладных сервисов;
\item \emph{межсетевые экраны, IDS/IPS, VPN-средства} — обеспечение безопасности при объединении сетей.
\end{itemize}
Маршрутизация обеспечивает доставку пакетов между сетями с различной адресацией. Инкапсуляция и туннелирование позволяют передавать пакеты одной сети через сети с другими протоколами. NAT решает проблему нехватки глобальных адресов и скрывает внутреннюю структуру локальной сети. Шлюзы и прокси-серверы обеспечивают взаимодействие прикладных программ, использующих разные протоколы. Средства безопасности защищают сеть от несанкционированного доступа и атак.
\textit{(Источник: Степанов, ч. III, гл. 21–26, с. 451–482.)}

\paragraph{Виртуальные частные сети.}
Виртуальная частная сеть (VPN) представляет собой логическую сеть, построенную поверх другой сети, обычно Интернета. VPN обеспечивает защищённую передачу данных за счёт шифрования и туннелирования. Она позволяет удалённым пользователям и филиалам организации подключаться к корпоративной сети через публичные каналы, сохраняя конфиденциальность и целостность информации. VPN является важным программно-техническим средством объединения территориально распределённых сетей.
\textit{(Источник: Степанов, ч. III, гл. 25–26, с. 472–482.)}

\paragraph{Управление сетью и мониторинг.}
Для объединения и эксплуатации различных сетей необходимы средства управления: протоколы управления (например, SNMP), системы мониторинга, анализаторы трафика, средства диагностики и конфигурирования. Они позволяют контролировать состояние оборудования, выявлять неисправности, анализировать загрузку каналов и обеспечивать согласованную работу разнородных сетевых компонентов.
\textit{(Источник: Степанов, ч. III, гл. 21–26, с. 451–482.)}

\subsubsection{Краткое резюме подраздела}
Локальные сети охватывают ограниченную территорию, характеризуются высокой скоростью и простотой управления; глобальные сети имеют большую протяжённость, используют арендованные каналы и сложную маршрутизацию. Техническими средствами объединения сетей являются повторители, концентраторы, мосты, коммутаторы, маршрутизаторы и шлюзы, работающие на разных уровнях модели OSI, а также линии связи и кабельная инфраструктура. Программную основу объединения составляют стеки протоколов, сетевые операционные системы, службы, маршрутизация, туннелирование, NAT, VPN и средства безопасности. Совместное применение технических и программных средств обеспечивает построение масштабируемых, надёжных и совместимых информационно-вычислительных сетей.
\textit{(Обобщение по: Степанов, ч. III, гл. 19–26, с. 432–482.)}

\subsection{Методы и средства передачи данных в ИВС, протоколы передачи данных}

\noindent\textit{Основные источники: Олифер В.Г., Олифер Н.А. Компьютерные сети. Принципы, технологии, протоколы. 2-е изд. Гл. 2, 4–5, 7, 11–14.}

\subsubsection{Методы передачи данных}

\paragraph{Линии связи и их характеристики.}
Передача данных в ИВС осуществляется по линиям связи, которые включают физическую среду, аппаратуру передачи данных и промежуточную аппаратуру. Физической средой могут быть медные провода, витая пара, коаксиальный кабель, волоконно-оптический кабель, радиоканалы, спутниковые и инфракрасные каналы. Линии связи делятся на проводные (воздушные), кабельные (медные и волоконно-оптические) и радиоканалы наземной и спутниковой связи. Основными характеристиками линий являются полоса пропускания, затухание, помехоустойчивость, волновое сопротивление, задержка распространения сигнала и пропускная способность. Пропускная способность зависит от полосы пропускания, уровня помех и способа кодирования. Теоретические пределы определяются формулой Шеннона
\[
C = F \log_2\left(1 + \frac{P_c}{P_{\text{ш}}}\right),
\]
где \(C\) — максимальная пропускная способность, \(F\) — ширина полосы пропускания, \(P_c\) — мощность сигнала, \(P_{\text{ш}}\) — мощность шума, а также формулой Найквиста
\[
C = 2F \log_2 M,
\]
где \(M\) — количество различимых состояний информационного параметра.
\textit{(Источник: Олифер, гл. 4.)}

\paragraph{Аналоговая модуляция.}
Аналоговая модуляция применяется для передачи дискретных данных по каналам с узкой полосой частот, типичным представителем которых является канал тональной частоты с полосой 3100 Гц. При амплитудной модуляции (АМ) для логической единицы выбирается один уровень амплитуды несущей, для нуля — другой. При частотной модуляции (ЧМ) значения 0 и 1 передаются синусоидами с различной частотой. При фазовой модуляции (ФМ) значениям данных соответствуют сигналы одинаковой частоты, но различной фазы. В скоростных модемах применяются комбинированные методы, прежде всего квадратурная амплитудная модуляция (QAM), сочетающая фазовую модуляцию с 8 значениями сдвига фазы и амплитудную модуляцию с 4 уровнями амплитуды. Устройство, выполняющее модуляцию и демодуляцию, называется модемом.
\textit{(Источник: Олифер, гл. 5.)}

\paragraph{Цифровое кодирование.}
При цифровом кодировании дискретной информации применяются потенциальные и импульсные коды. В потенциальных кодах для представления логических единиц и нулей используется только значение потенциала сигнала, а его перепады, формирующие законченные импульсы, во внимание не принимаются. Импульсные коды представляют двоичные данные либо импульсами определенной полярности, либо частью импульса — перепадом потенциала определенного направления. К основным потенциальным кодам относятся NRZ (без возвращения к нулю), AMI (биполярное кодирование с альтернативной инверсией), NRZI (потенциальный код с инверсией при единице) и 2B1Q (потенциальный код с четырьмя уровнями сигнала). К импульсным относится биполярный импульсный код и манчестерский код. Для улучшения свойств потенциальных кодов применяются избыточные коды (4B/5B, 8B/6T) и скрэмблирование (B8ZS, HDB3). Метод 4B/5B заменяет исходные символы длиной 4 бита символами длиной 5 бит, что позволяет устранить длинные последовательности нулей и придать коду свойство самосинхронизации. Скрэмблирование заключается в побитном вычислении результирующего кода на основании битов исходного кода и полученных ранее битов результирующего кода.
\textit{(Источник: Олифер, гл. 5.)}

\paragraph{Дискретная модуляция аналоговых сигналов.}
Для передачи аналоговых сигналов в цифровой форме применяется дискретная модуляция, основанная на дискретизации непрерывных процессов по амплитуде и по времени. Типичным примером является импульсно-кодовая модуляция (ИКМ). Амплитуда исходной непрерывной функции измеряется с заданным периодом, затем каждый замер представляется двоичным числом определенной разрядности. Устройство, выполняющее эту функцию, называется аналого-цифровым преобразователем (АЦП). На приемной стороне цифро-аналоговый преобразователь (ЦАП) восстанавливает непрерывный сигнал. В соответствии с теоремой Котельникова частота дискретизации должна быть не менее чем в два раза выше самой высокой гармоники спектра исходной функции. Для качественной передачи голоса используется частота квантования 8000 Гц и 7 или 8 бит на замер, что дает скорость 56 или 64 кбит/с. Применяются также дифференциальные методы ADPCM и LPC, позволяющие снизить скорость передачи голоса.
\textit{(Источник: Олифер, гл. 5.)}

\paragraph{Мультиплексирование.}
Для совместного использования одной линии связи несколькими потоками применяются частотное мультиплексирование (FDM), временное мультиплексирование (TDM) и спектральное мультиплексирование по длине волны (WDM/DWDM). В сетях FDM каждому абонентскому каналу выделяется полоса частот шириной около 4 кГц; существует иерархия каналов FDM, в которой 12 абонентских каналов образуют базовую группу, 5 базовых групп — супергруппу, 10 супергрупп — главную группу. В сетях TDM каждому каналу выделяется определенная доля времени (тайм-слот) в цикле работы оборудования; цикл обычно равен 125 мкс, что соответствует периоду следования замеров голоса. Мультиплексирование по длине волны WDM/DWDM использует тот же принцип частотного разделения, но информационным сигналом является свет. В магистральном канале обычно мультиплексируется несколько спектральных каналов — до 16, 32, 40, 80 или 160.
\textit{(Источник: Олифер, гл. 6.)}

\paragraph{Методы коммутации.}
В ИВС применяются коммутация каналов, коммутация сообщений и коммутация пакетов. В сетях с коммутацией каналов перед передачей данных всегда выполняется процедура установления соединения, после которой создается составной канал с фиксированной пропускной способностью. Коммутация сообщений предназначена для организации взаимодействия пользователей в автономном режиме; сообщение передается через несколько транзитных компьютеров, где оно целиком буферизуется на диске. Сети с коммутацией пакетов были специально разработаны для эффективной передачи пульсирующего компьютерного трафика. Буферизация пакетов разных абонентов в коммутаторах позволяет сгладить неравномерности интенсивности трафика каждого абонента и равномерно загрузить каналы связи между коммутаторами. В сетях с коммутацией пакетов применяются два класса механизмов продвижения: дейтаграммная передача и виртуальные каналы. При дейтаграммной передаче все передаваемые пакеты обрабатываются независимо друг от друга, пакет за пакетом; выбор следующего узла происходит только на основании адреса назначения, содержащегося в заголовке пакета. При виртуальных каналах для всех пакетов потока прокладывается единый путь через сеть; решение о продвижении принимается на основании локального идентификатора виртуального канала. Виртуальные каналы могут быть коммутируемыми (SVC) и постоянными (PVC).
\textit{(Источник: Олифер, гл. 2.)}

\paragraph{Обеспечение достоверности передачи.}
Для обнаружения и коррекции ошибок применяются контрольные суммы, циклический избыточный код (CRC), контроль по паритету, квитанции, тайм-ауты, повторная передача и алгоритм скользящего окна. Метод CRC основан на рассмотрении исходных данных в виде одного многоразрядного двоичного числа и вычислении остатка от деления этого числа на известный делитель. Метод CRC обнаруживает все одиночные ошибки, двойные ошибки и ошибки в нечетном числе битов. Для восстановления искаженных и потерянных кадров используются методы с простоями и скользящего окна. В методе с простоями источник, пославший кадр, ожидает получения квитанции от приемника и только после этого посылает следующий кадр. В методе скользящего окна источнику разрешается передать некоторое количество кадров в непрерывном режиме без получения положительных ответных квитанций; количество таких кадров называется размером окна.
\textit{(Источник: Олифер, гл. 5.)}

\paragraph{Компрессия данных.}
Для сокращения времени передачи данных применяется компрессия, или сжатие, данных. Различают статическую и динамическую компрессию. К общим алгоритмам компрессии относятся десятичная упаковка, относительное кодирование, символьное подавление, коды переменной длины (в том числе алгоритм Хаффмана). Десятичная упаковка использует простое двоичное кодирование десятичных цифр вместо кода ASCII. Относительное кодирование передает только отклонения между последовательными цифрами вместе с известным опорным значением. Символьное подавление заменяет последовательность из трех или более одинаковых байтов специальной трехбайтовой последовательностью. Коды переменной длины заменяют часто встречающиеся символы кодами меньшей длины, а редко встречающиеся — кодами большей длины. Динамическая компрессия выполняется в процессе передачи данных; реальный коэффициент компрессии зависит от типа передаваемых данных и применяемого алгоритма.
\textit{(Источник: Олифер, гл. 5.)}

\subsubsection{Средства передачи данных}

\paragraph{Кабельные системы.}
В компьютерных сетях применяются три основных типа кабеля: кабели на основе скрученных пар медных проводов (витая пара), коаксиальные кабели с медной жилой, а также волоконно-оптические кабели. Витая пара делится на неэкранированную (UTP) и экранированную (STP). Кабели UTP проще в изготовлении и монтаже, зато кабели STP обеспечивают более высокий уровень защищенности. Кабели на основе неэкранированной витой пары делятся на категории: категория 3 поддерживает частоты до 16 МГц, категория 4 — до 20 МГц, категория 5 — до 100 МГц, категория 5е — до 100 МГц с улучшенными характеристиками, категория 6 — до 250 МГц, категория 7 — до 600 МГц. Волоконно-оптические кабели делятся на многомодовые и одномодовые. Многомодовые кабели применяются для передачи данных на небольшие расстояния со скоростью до 1 Гбит/с. Одномодовые кабели являются основой сверхскоростных дальних линий связи, обеспечивающих передачу информации со скоростями в сотни и тысячи гигабит в секунду на расстояние в несколько сот километров.
\textit{(Источник: Олифер, гл. 4, 9.)}

\paragraph{Сетевые адаптеры.}
Сетевой адаптер (Network Interface Card, NIC) вместе со своим драйвером реализует второй, канальный уровень модели открытых систем в конечном узле сети — компьютере. Сетевой адаптер совместно с драйвером выполняет две операции: передачу и прием кадра. Передача кадра включает прием кадра данных LLC через межуровневый интерфейс, оформление кадра данных MAC-уровня, формирование символов кодов при использовании избыточных кодов, скрэмблирование кодов, выдачу сигналов в кабель в соответствии с принятым линейным кодом. Прием кадра включает прием из кабеля сигналов, выделение сигналов на фоне шума, дескрэмблирование, проверку контрольной суммы кадра, извлечение кадра LLC и передачу его протоколу LLC. Сетевые адаптеры делятся на адаптеры для клиентских компьютеров и адаптеры для серверов. В адаптерах для клиентских компьютеров значительная часть работы перекладывается на драйвер, тем самым адаптер оказывается проще и дешевле. Адаптеры, предназначенные для серверов, обычно снабжаются собственными процессорами, которые самостоятельно выполняют большую часть работы по передаче кадров из оперативной памяти в сеть и в обратном направлении.
\textit{(Источник: Олифер, гл. 9.)}

\paragraph{Модемы.}
Модем — устройство, которое выполняет функции модуляции несущей синусоиды на передающей стороне и демодуляции на приемной стороне. Модемы для выделенных каналов работают на основе методов аналоговой модуляции сигнала. Протоколы и стандарты модемов определены в рекомендациях CCITT серии V и делятся на три группы: стандарты, определяющие скорость передачи данных и метод кодирования; стандарты исправления ошибок; стандарты сжатия данных. Модемы, поддерживающие только асинхронный режим, обычно обеспечивают низкую скорость передачи данных — до 1200 бит/с. Модемы, работающие только в синхронном режиме, могут подключаться только к 4-проводному окончанию. Модемы, работающие в асинхронном и синхронном режимах, являются наиболее универсальными устройствами. К основным стандартам модемов для коммутируемых каналов относятся V.21, V.22, V.22 bis, V.26 ter, V.32, V.32 bis, V.34, V.34+, V.90, V.92. Стандарт V.34 разрабатывался для передачи информации по каналам практически любого качества; его особенностью являются процедуры динамической адаптации к изменениям характеристик канала во время обмена информацией. Стандарт V.90 обеспечивает асимметричный обмен данными: со скоростью 56 кбит/с из сети и со скоростью 33,6 кбит/с в сеть. Стандарт V.42 определяет протокол коррекции ошибок LAP-M, а стандарт V.42 bis — метод сжатия данных при асинхронном интерфейсе.
\textit{(Источник: Олифер, гл. 17.)}

\paragraph{Устройства DSU/CSU и терминальные адаптеры ISDN.}
Связь компьютера или маршрутизатора с цифровой выделенной линией PDH/SDH осуществляется с помощью пары устройств, обычно выполненных в одном корпусе или совмещенных с маршрутизатором. Этими устройствами являются устройство обслуживания данных (DSU) и устройство обслуживания канала (CSU). DSU преобразует сигналы, поступающие от DTE, выполняет синхронизацию, формирует кадры каналов E1–E4 или STM-N, усиливает сигнал и осуществляет выравнивание загрузки канала. Устройство CSU выполняет более узкие функции, в основном создает оптимальные условия передачи в линии. Для подключения к коммутируемому каналу ISDN используются терминальные адаптеры (Terminal Adapter, TA). Эти устройства преобразуют сигналы интерфейсов DTE-DCE, например интерфейса V.35, в сигналы интерфейса точки S/T стандарта ISDN.
\textit{(Источник: Олифер, гл. 17.)}

\paragraph{Интерфейсы DTE-DCE.}
Для подключения устройств DCE к аппаратуре, вырабатывающей данные для глобальной сети, то есть устройствам DTE, существует несколько стандартных интерфейсов, которые представляют собой стандарты физического уровня. К этим стандартам относятся стандарты серии V CCITT, а также стандарты EIA серии RS. Интерфейс RS-232C/V.24 является наиболее популярным низкоскоростным интерфейсом. Первоначально он был разработан для передачи данных между компьютером и модемом со скоростью не выше 9600 бит/с на расстояние до 15 м. Позднее практические реализации этого интерфейса стали работать и на более высоких скоростях — до 115 200 бит/с. Интерфейс поддерживает как асинхронный, так и синхронный режимы работы. Интерфейс RS-449/V.10/V.11 поддерживает более высокую скорость обмена данными и большую удаленность DCE от DTE. Интерфейс V.35 был разработан для подключения синхронных модемов; он обеспечивает только синхронный режим обмена между DTE и DCE на скорости до 168 кбит/с. Интерфейс HSSI (High-Speed Serial Interface) разработан для подключения к устройствам DCE, работающим на высокоскоростные каналы, такие как каналы T3 (45 Мбит/с) и SONET OC-1 (52 Мбит/с). Интерфейс работает в синхронном режиме и поддерживает передачу данных в диапазоне скоростей от 300 кбит/с до 52 Мбит/с.
\textit{(Источник: Олифер, гл. 17.)}

\subsubsection{Протоколы передачи данных}

\paragraph{Понятие протокола и интерфейса.}
Протокол — это формально определенная процедура взаимодействия объектов, а интерфейс — набор правил взаимодействия между уровнями. Протоколы реализуются не только компьютерами, но и другими сетевыми устройствами — концентраторами, мостами, коммутаторами, маршрутизаторами. В общем случае связь компьютеров в сети осуществляется не напрямую, а через различные коммуникационные устройства. В зависимости от типа устройства в нем должны быть встроены средства, реализующие тот или иной набор протоколов. Протоколы реализуются в виде программных модулей и выполняются на конечных узлах-компьютерах, а также на промежуточных узлах-маршрутизаторах.
\textit{(Источник: Олифер, гл. 3.)}

\paragraph{Многоуровневая модель OSI.}
Модель взаимодействия открытых систем (Open System Interconnection, OSI) определяет различные уровни взаимодействия систем в сетях с коммутацией пакетов, дает им стандартные имена и указывает, какие функции должен выполнять каждый уровень. В модели OSI средства взаимодействия делятся на семь уровней: прикладной, представительный, сеансовый, транспортный, сетевой, канальный и физический. Каждый уровень имеет дело с одним определенным аспектом взаимодействия сетевых устройств. Физический уровень имеет дело с передачей битов по физическим каналам связи. Канальный уровень обеспечивает передачу пакетов данных, поступающих от протоколов верхних уровней, узлу назначения, адрес которого также указывает протокол верхнего уровня; протоколы канального уровня оформляют переданные им пакеты в кадры собственного формата, помещая указанный адрес назначения в одно из полей такого кадра, а также сопровождая кадр контрольной суммой. Сетевой уровень служит для образования единой транспортной системы, объединяющей несколько сетей, причем эти сети могут использовать совершенно разные принципы передачи сообщений между конечными узлами и обладать произвольной структурой связей. Транспортный уровень обеспечивает приложениям или верхним уровням стека — прикладному и сеансовому — передачу данных с той степенью надежности, которая им требуется. Сеансовый уровень обеспечивает управление взаимодействием: фиксирует, какая из сторон является активной в настоящий момент, предоставляет средства синхронизации. Представительный уровень имеет дело с формой представления передаваемой по сети информации, не меняя при этом ее содержания. Прикладной уровень — самый верхний уровень, являющийся, по сути, набором разнообразных протоколов, с помощью которых пользователи сети получают доступ к разделяемым ресурсам, таким как файлы, принтеры или гипертекстовые web-страницы.
\textit{(Источник: Олифер, гл. 3.)}

\paragraph{Стек TCP/IP.}
Стек TCP/IP был разработан по инициативе Министерства обороны США более 20 лет назад для связи экспериментальной сети ARPAnet с другими сетями как набор общих протоколов для разнородной вычислительной среды. Стек TCP/IP имеет 4 уровня: прикладной, транспортный, уровень межсетевого взаимодействия и уровень сетевых интерфейсов. Прикладной уровень объединяет службы, предоставляемые системой пользовательским приложениям: протокол копирования файлов FTP, протокол эмуляции терминала telnet, простой протокол передачи электронной почты SMTP, протокол передачи гипертекстовой информации HTTP и многие другие. На транспортном уровне стека TCP/IP функционируют протоколы TCP и UDP. Протокол управления передачей TCP решает задачу обеспечения надежной информационной связи между двумя конечными узлами. Дейтаграммный протокол UDP используется как экономичное средство связи уровня межсетевого взаимодействия с прикладным уровнем. Уровень межсетевого взаимодействия реализует концепцию коммутации пакетов в режиме без установления соединений. Основными протоколами этого уровня являются дейтаграммный протокол IP и протоколы маршрутизации (RIP, OSPF, BGP и др.). Вспомогательную роль выполняют протокол межсетевых управляющих сообщений ICMP и протокол группового управления IGMP. Протоколы уровня сетевых интерфейсов обеспечивают интеграцию в составную сеть других сетей. Этот уровень не регламентируется, но поддерживает все популярные стандарты физического и канального уровней: для локальных сетей — Ethernet, Token Ring, FDDI и т. д., для глобальных сетей — X.25, frame relay, PPP, ISDN и т. д. К этому уровню часто относят и протокол разрешения адресов ARP.
\textit{(Источник: Олифер, гл. 11.)}

\paragraph{Протоколы канального уровня.}
В локальных сетях применяются протоколы Ethernet, Token Ring, FDDI. Ethernet использует метод доступа CSMA/CD, MAC-адресацию и кадры Ethernet. Кадр Ethernet содержит адреса назначения и источника, поле длины/типа, поле данных и контрольную сумму. Минимальная длина кадра с учетом служебных полей составляет 64 байта, максимальная — 1518 байт без преамбулы. Token Ring использует маркерный метод доступа и логическую топологию кольца. Станция может передавать данные только после получения маркера. Кадры Token Ring содержат начальный и конечный ограничители, поле управления доступом, адреса, данные, контрольную сумму и поле статуса. FDDI — технология локальных сетей со скоростью 100 Мбит/с, разработанная ANSI. Она использует двойное волоконно-оптическое кольцо, маркерный метод доступа и обеспечивает высокую отказоустойчивость. Для глобальных связей типа «точка-точка» применяются протоколы SLIP, HDLC и PPP. Протокол SLIP выполняет единственную функцию — позволяет в потоке битов, которые поступают по выделенному или коммутируемому каналу, распознать начало и конец IP-пакета. Протокол PPP разработан группой IETF как часть стека TCP/IP для передачи кадров информации по последовательным глобальным каналам связи взамен устаревшего протокола SLIP. Основное отличие PPP от других протоколов канального уровня состоит в том, что он добивается согласованной работы различных устройств с помощью переговорной процедуры, во время которой передаются различные параметры, такие как качество линии, протокол аутентификации и инкапсулируемые протоколы сетевого уровня. Протокол LAP-B используется в сетях X.25, LAP-D — в сетях ISDN, LAP-F — в сетях frame relay.
\textit{(Источник: Олифер, гл. 5, 7, 17.)}

\paragraph{Протокол IP.}
Протокол IP относится к протоколам без установления соединения. Перед IP не ставится задача надежной доставки сообщений от отправителя к получателю. Протокол IP обрабатывает каждый IP-пакет как независимую единицу, не имеющую связи ни с какими другими IP-пакетами. В протоколе IP нет механизмов, обычно применяемых для увеличения достоверности конечных данных: отсутствует квитирование, нет процедуры упорядочивания, повторных передач или других подобных функций. Важной особенностью протокола IP является его способность выполнять динамическую фрагментацию пакетов при передаче их между сетями с различными максимально допустимыми значениями длины поля данных кадров (MTU). IP-пакет состоит из заголовка и поля данных. Заголовок обычно имеет длину 20 байт и содержит информацию о сетевых адресах отправителя и получателя, о параметрах фрагментации, о времени жизни пакета, о контрольной сумме и некоторых других. В поле данных IP-пакета находятся сообщения более высокого уровня, например TCP или UDP. IP-адрес имеет длину 4 байта и состоит из номера сети и номера узла. Для определения границы, отделяющей номер сети от номера узла, реализуются два подхода: первый основан на понятии класса адреса, второй — на использовании масок. Класс адреса определяется значениями нескольких первых битов адреса. В адресах класса A под номер сети отводится один байт, а остальные три байта — под номер узла. Для небольших сетей больше подходят адреса класса C, в которых номер сети занимает три байта, а для нумерации узлов может быть использован только один байт. Другой способ определения, какая часть адреса является номером сети, а какая номером узла, основан на использовании маски. Маска — это число, которое используется в паре с IP-адресом; двоичная запись маски содержит единицы в тех разрядах, которые в IP-адресе должны интерпретироваться как номер сети.
\textit{(Источник: Олифер, гл. 12–13.)}

\paragraph{Протоколы TCP и UDP.}
Протокол TCP работает непосредственно над протоколом IP и использует для транспортировки своих блоков данных этот дейтаграммный и поэтому потенциально ненадежный протокол. Надежность передачи данных протоколом TCP достигается за счет того, что он основан на установлении логических соединений между взаимодействующими процессами. До тех пор пока модули протокола TCP на обоих оконечных хостах продолжают функционировать корректно, а составная сеть не распалась на несвязные части, ошибки в передаче данных на уровне протокола IP не будут влиять на правильность обмена хостов сегментами. Информация, поступающая к протоколу TCP в рамках логического соединения от протоколов более высокого уровня, рассматривается протоколом TCP как неструктурированный поток байтов. Поступающие данные буферизуются средствами TCP. Для передачи на сетевой уровень из буфера «вырезается» некоторая непрерывная часть данных, которая называется сегментом. В заголовок сегмента протокол TCP помещает номер порта назначения и номер порта отправителя. Каждая сторона соединения посылает противоположной стороне следующие параметры: максимальный размер сегмента, который она готова принимать; максимальный объем данных, которые она разрешает другой стороне передавать в свою сторону; начальный порядковый номер байта, с которого она начинает отсчет потока данных в рамках данного соединения. В рамках установленного соединения правильность передачи каждого сегмента должна подтверждаться квитанцией получателя. В протоколе TCP используется алгоритм скользящего окна. Особенность использования алгоритма скользящего окна в протоколе TCP состоит в том, что, хотя единицей передаваемых данных является сегмент, окно определено на множестве нумерованных байтов неструктурированного потока данных. UDP является простейшим дейтаграммным протоколом, который используется в том случае, когда задача надежного обмена данными либо вообще не ставится, либо решается средствами более высокого уровня — системными прикладными службами или пользовательскими приложениями.
\textit{(Источник: Олифер, гл. 13.)}

\paragraph{Служебные протоколы и маршрутизация.}
Для определения локального адреса по IP-адресу используется протокол разрешения адресов ARP. Протокол разрешения адресов реализуется различным образом в зависимости от того, какой протокол канального уровня работает в данной сети — протокол локальной сети (Ethernet, Token Ring, FDDI) с возможностью широковещательного доступа одновременно ко всем узлам сети или же какой-либо из протоколов глобальной сети (X.25, frame relay), которые, как правило, не поддерживают широковещательный доступ. Протокол межсетевых управляющих сообщений ICMP предназначен для обмена информацией об ошибках между маршрутизаторами сети и узлом-источником пакета. С помощью специальных пакетов ICMP сообщает о невозможности доставки пакета, о превышении времени жизни или продолжительности сборки пакета из фрагментов, об аномальных величинах параметров, об изменении маршрута пересылки и типа обслуживания, о состоянии системы и т. п. Для автоматического построения таблиц маршрутизации маршрутизаторы обмениваются информацией о топологии составной сети в соответствии со специальным служебным протоколом. Протоколы этого типа называются протоколами маршрутизации. Наиболее распространенным протоколом, основанным на дистанционно-векторном алгоритме, является протокол RIP. Протокол OSPF является достаточно современной реализацией алгоритма состояния связей и обладает многими особенностями, ориентированными на применение в больших гетерогенных сетях. Внешний протокол BGP отвечает за выбор маршрута как последовательности автономных систем.
\textit{(Источник: Олифер, гл. 12–14.)}

\paragraph{Прикладные протоколы.}
К прикладным протоколам стека TCP/IP относятся протокол копирования файлов FTP, протокол эмуляции терминала telnet, простой протокол передачи электронной почты SMTP, протокол передачи гипертекстовой информации HTTP, система доменных имен DNS, протокол динамического конфигурирования узлов DHCP, протокол управления сетями SNMP. Протокол FTP предназначен для обмена файлами между компьютерами сети. Протокол telnet позволяет пользователю удаленно работать с компьютером таким же способом, как если бы он управлял им с помощью локально подключенного терминала. Протокол SMTP используется для передачи электронной почты. Протокол HTTP является основным протоколом службы WWW и обеспечивает передачу гипертекстовой информации. Система доменных имен DNS устанавливает соответствие между доменными именами и IP-адресами. Протокол DHCP автоматизирует процесс назначения IP-адресов. Протокол SNMP используется для управления сетевыми устройствами и получения от них информации об их статусе, производительности и других характеристиках.
\textit{(Источник: Олифер, гл. 11–12, 20.)}

\begin{table}[h]
\centering
\caption{Примеры протоколов по уровням модели OSI и стека TCP/IP}
\label{tab:protocols-osi-tcpip}
\begin{tabular}{p{3.2cm}p{4cm}p{6.5cm}}
\toprule
\textbf{Уровень OSI} & \textbf{Уровень TCP/IP} & \textbf{Примеры протоколов и технологий} \\
\midrule
Прикладной, представительный, сеансовый & Прикладной & HTTP, FTP, SMTP, DNS, DHCP, SNMP, telnet \\
Транспортный & Транспортный & TCP, UDP \\
Сетевой & Межсетевого взаимодействия & IP, ICMP, ARP, RIP, OSPF, BGP \\
Канальный & Сетевых интерфейсов & Ethernet, Token Ring, FDDI, PPP, SLIP, HDLC, LAP-B, LAP-D, LAP-F \\
Физический & Сетевых интерфейсов & 10Base-T, 100Base-TX, 100Base-FX, V.35, RS-232C, SDH, DWDM \\
\bottomrule
\end{tabular}
\end{table}

\subsubsection{Краткое резюме подраздела}
Методы передачи данных в ИВС включают выбор среды, кодирование и модуляцию, мультиплексирование, коммутацию, обеспечение достоверности и сжатие. Средства передачи данных включают кабельные системы, сетевые адаптеры, модемы, DSU/CSU, терминальные адаптеры, повторители, концентраторы, коммутаторы и маршрутизаторы. Взаимодействие устройств описывается многоуровневыми наборами протоколов. Модель OSI определяет семь уровней взаимодействия; стек TCP/IP включает прикладной, транспортный, межсетевой уровни и уровень сетевых интерфейсов. Ключевую роль играют протоколы Ethernet, IP, TCP, UDP, ICMP, ARP, DNS, DHCP и протоколы маршрутизации RIP, OSPF, BGP. Совместное применение методов, средств и протоколов обеспечивает построение масштабируемых, надёжных и совместимых информационно-вычислительных сетей.
\textit{(Обобщение по: Олифер, гл. 2–7, 11–14, 17, 20.)}

\subsection{Особенности архитектуры локальных сетей (Ethernet, Token Ring, FDDI)}

\noindent\textit{Основные источники: Олифер В.Г., Олифер Н.А. Компьютерные сети. Принципы, технологии, протоколы. 2-е изд. Гл. 7–8.}

\subsubsection{Общие принципы построения локальных сетей}

\paragraph{Разделяемая среда и типовые топологии.}
Классические технологии локальных сетей — Ethernet, Token Ring, FDDI — строятся на основе разделяемой среды передачи данных. Разделяемая среда представляет собой физический канал, к которому подключены все станции сети; в каждый момент времени передавать данные может только одна станция. Для доступа к среде используются специальные методы: случайный доступ с обнаружением коллизий (CSMA/CD) в Ethernet и маркерный доступ в Token Ring и FDDI. Типовыми топологиями локальных сетей являются общая шина, звезда, кольцо и иерархическая звезда (дерево). Физическая топология может отличаться от логической: например, в Token Ring физически станции подключаются к концентраторам по звездообразной схеме, а логически образуют кольцо. Выбор топологии и метода доступа определяет производительность, надёжность, стоимость и масштабируемость сети.
\textit{(Источник: Олифер, гл. 7.)}

\paragraph{Уровни MAC и LLC.}
В стандартах IEEE 802 канальный уровень разделён на два подуровня: управление логическим каналом (Logical Link Control, LLC) и управление доступом к среде (Media Access Control, MAC). Уровень MAC обеспечивает корректное совместное использование общей среды и реализует конкретный метод доступа (CSMA/CD, маркерный). Уровень LLC, определённый стандартом 802.2, предоставляет верхним уровням три типа процедур: LLC1 — без установления соединения и без подтверждения; LLC2 — с установлением соединения и подтверждением; LLC3 — без установления соединения, но с подтверждением. Протокол LLC не зависит от конкретной технологии и позволяет различным сетевым технологиям взаимодействовать с протоколами верхних уровней.
\textit{(Источник: Олифер, гл. 7.)}

\paragraph{Форматы кадров и адресация.}
Все технологии локальных сетей используют 6-байтовые MAC-адреса, уникальные для каждого сетевого интерфейса. MAC-адрес назначается производителем и содержит идентификатор организации-изготовителя (OUI) в трёх старших байтах. Первый бит старшего байта адреса назначения указывает на индивидуальный (unicast) или групповой (multicast) адрес; широковещательный адрес состоит из всех единиц. Кадры локальных сетей содержат адреса назначения и источника, служебные поля, поле данных и контрольную сумму. Форматы кадров Ethernet, Token Ring и FDDI различаются, но все они обеспечивают передачу кадров LLC верхним уровням.
\textit{(Источник: Олифер, гл. 7.)}

\subsubsection{Технология Ethernet}

\paragraph{История и семейство стандартов.}
Ethernet — самый распространённый стандарт локальных сетей. Первоначально разработан компаниями Xerox, Digital и Intel (стандарт Ethernet DIX), затем на его основе был принят стандарт IEEE 802.3. В семейство Ethernet входят классический Ethernet 10 Мбит/с, Fast Ethernet 100 Мбит/с, Gigabit Ethernet 1000 Мбит/с и 10 Gigabit Ethernet. Все они используют единый формат MAC-адресов и, как правило, один и тот же метод доступа CSMA/CD в полудуплексном режиме; в полнодуплексном режиме, который применяется при подключении к коммутаторам, коллизии отсутствуют.
\textit{(Источник: Олифер, гл. 7–8.)}

\paragraph{Метод доступа CSMA/CD.}
Метод CSMA/CD (Carrier Sense Multiple Access with Collision Detection) включает следующие этапы. Станция, желающая передавать, прослушивает среду (carrier sense). Если среда свободна, станция начинает передачу кадра. Во время передачи она продолжает прослушивать среду и сравнивает передаваемый сигнал с принимаемым. Если обнаруживается расхождение, фиксируется коллизия (collision detection). Станция прерывает передачу, посылает в сеть jam-последовательность из 32 бит и прекращает передачу. Затем она выжидает случайный интервал времени и повторяет попытку. Случайная пауза выбирается по алгоритму двоичной экспоненциальной задержки: после \(n\)-й коллизии интервал выбирается из диапазона \([0, 2^n-1]\) слотов отсрочки, где слот равен 512 битовым интервалам. После 10 попыток диапазон не увеличивается, после 16 попыток кадр отбрасывается. Метод CSMA/CD является вероятностным и не гарантирует станции получение доступа к среде; при высокой загрузке сети вероятность коллизий растёт, а полезная пропускная способность падает.
\textit{(Источник: Олифер, гл. 7.)}

\paragraph{Время двойного оборота и минимальная длина кадра.}
Для надёжного распознавания коллизий всеми станциями должно выполняться условие: время передачи кадра минимальной длины \(T_{\min}\) должно быть не меньше времени двойного оборота сигнала \(PDV\) между двумя наиболее удалёнными станциями сети. В стандарте Ethernet минимальная длина поля данных кадра составляет 46 байт, что вместе со служебными полями даёт минимальную длину кадра 64 байта (без преамбулы), или 576 бит. Время передачи такого кадра при скорости 10 Мбит/с равно 57,5 мкс. Это накладывает ограничения на максимальный диаметр сети и количество повторителей. В Fast Ethernet и Gigabit Ethernet с увеличением скорости время передачи кадра минимальной длины уменьшается, поэтому для сохранения допустимого диаметра сети в Gigabit Ethernet минимальный размер кадра был увеличен до 512 байт (с использованием поля расширения) и разрешён режим пакетного режима (burst mode).
\textit{(Источник: Олифер, гл. 7–8.)}

\paragraph{Форматы кадров Ethernet.}
В сетях Ethernet используются четыре основных формата кадров: кадр 802.3/LLC (802.3/802.2), кадр Raw 802.3 (Novell 802.3), кадр Ethernet DIX (Ethernet II) и кадр Ethernet SNAP. Кадр 802.3/LLC включает заголовок MAC (адреса, поле длины) и вложенный кадр LLC с полями DSAP, SSAP и управляющим полем. Кадр Raw 802.3 не содержит заголовка LLC; его поле данных начинается непосредственно с пакета IPX. Кадр Ethernet DIX использует поле Type (2 байта) вместо поля длины для указания протокола верхнего уровня. Кадр Ethernet SNAP является расширением кадра 802.3/LLC и содержит дополнительный заголовок SNAP с полями OUI и Type, что обеспечивает совместимость с кодами протоколов Ethernet II. Сетевые адаптеры и коммутаторы автоматически распознают тип кадра.
\textit{(Источник: Олифер, гл. 7.)}

\paragraph{Физический уровень Ethernet.}
Классический Ethernet 10 Мбит/с имеет следующие спецификации физического уровня: 10Base-5 (толстый коаксиальный кабель, максимальная длина сегмента 500 м), 10Base-2 (тонкий коаксиальный кабель, 185 м), 10Base-T (неэкранированная витая пара, звездообразная топология с концентратором, до 100 м), 10Base-F (волоконно-оптический кабель, варианты FOIRL, 10Base-FL, 10Base-FB). Для коаксиальных сетей действует правило «5-4-3»: не более 5 сегментов, 4 повторителей и 3 нагруженных сегментов. Для витой пары и оптоволокна используется правило «4 хабов»: между любыми двумя станциями не должно быть более 4 концентраторов. Fast Ethernet использует спецификации 100Base-TX (две пары UTP категории 5 или STP Type 1), 100Base-T4 (четыре пары UTP категории 3) и 100Base-FX (многомодовое оптоволокно). Gigabit Ethernet работает по многомодовому и одномодовому оптоволокну, твинаксиальному кабелю и витой паре категории 5. В коммутируемых сетях Ethernet широко применяется полнодуплексный режим, при котором передача и приём ведутся одновременно, а коллизии отсутствуют.
\textit{(Источник: Олифер, гл. 7–8.)}

\paragraph{Производительность Ethernet.}
Максимальная пропускная способность сегмента Ethernet в кадрах в секунду достигается при передаче кадров минимальной длины и составляет 14 880 кадр/с для 10 Мбит/с. При этом полезная пропускная способность равна примерно 5,48 Мбит/с. При использовании кадров максимальной длины (1518 байт) максимальная пропускная способность составляет 813 кадр/с, а полезная — около 9,76 Мбит/с. Коэффициент использования сети при отсутствии коллизий и ожидания доступа может достигать 0,96. На практике при загрузке сегмента выше 30–50 % резко возрастает интенсивность коллизий и задержки, поэтому рекомендуется не превышать уровень загрузки 30 %.
\textit{(Источник: Олифер, гл. 7.)}

\subsubsection{Технология Token Ring}

\paragraph{Общая характеристика.}
Token Ring — технология локальных сетей, разработанная компанией IBM и стандартизованная IEEE 802.5. Она использует маркерный метод доступа и логическую топологию кольца. Физически станции подключаются к концентраторам MSAU (Multi-Station Access Unit) по топологии звезда, а концентраторы объединяются в кольцо через порты Ring In и Ring Out. Скорости работы — 4 и 16 Мбит/с. Смешение станций с разными скоростями в одном кольце не допускается. В сети назначается активный монитор — одна из станций, которая контролирует наличие маркера, выполняет ресинхронизацию сигналов и следит за временем оборота маркера.
\textit{(Источник: Олифер, гл. 7.)}

\paragraph{Маркерный метод доступа.}
В сети Token Ring по кольцу циркулирует специальный кадр — маркер. Станция, желающая передавать, должна дождаться маркера, захватить его и затем передать свои кадры данных. После завершения передачи станция передаёт маркер следующей станции. В сетях 4 Мбит/с станция удерживает маркер до передачи всех своих кадров, после чего передаёт маркер далее. В сетях 16 Мбит/с используется алгоритм раннего освобождения маркера (Early Token Release): станция передаёт маркер следующей станции сразу после окончания передачи последнего бита кадра, не дожидаясь возвращения кадра по кольцу. Это повышает эффективность использования кольца. Время удержания маркера ограничено; станция может передать один или несколько кадров в зависимости от размера кадров и величины времени удержания. Поддерживаются приоритеты кадров от 0 до 7; станция может захватить маркер только в том случае, если приоритет её кадра не ниже приоритета маркера.
\textit{(Источник: Олифер, гл. 7.)}

\paragraph{Активный монитор и отказоустойчивость.}
Одна из станций выполняет роль активного монитора. Она выбирается во время инициализации кольца как станция с максимальным MAC-адресом. Активный монитор генерирует кадр своего присутствия каждые 3 секунды; если этот кадр не появляется в сети более 7 секунд, остальные станции начинают процедуру выборов нового активного монитора. Активный монитор также удаляет из кольца кадры, которые не были удалены станцией-отправителем, и следит за временем оборота маркера. Если активный монитор не получает маркер в течение длительного времени, он порождает новый маркер. Для диагностики неисправностей используется процедура beaconing: станция, обнаружившая проблему, посылает кадры Beacon; последовательно определяются границы неисправного участка.
\textit{(Источник: Олифер, гл. 7.)}

\paragraph{Форматы кадров Token Ring.}
В Token Ring используются три формата кадров: маркер, кадр данных и прерывающая последовательность. Маркер состоит из полей начального ограничителя (SD), управления доступом (AC) и конечного ограничителя (ED). Кадр данных включает поля SD, управления кадром (FC), адреса назначения (DA), адреса источника (SA), данных, контрольной суммы (FCS), ED и статуса кадра (FS). Поле FS содержит биты распознавания адреса (A) и копирования кадра (C), которые дублируются для надёжности. Кадры MAC-уровня используются для управления кольцом: тест дублирования адреса (DAT), кадр присутствия активного монитора (AMP), кадр присутствия резервного монитора (SMP), маркер заявки (CT), сигнал (BCN) и очистка (PRC). Прерывающая последовательность состоит из начального и конечного ограничителей и сигнализирует об отмене текущей передачи.
\textit{(Источник: Олифер, гл. 7.)}

\paragraph{Физический уровень Token Ring.}
Сеть Token Ring строится на основе концентраторов MSAU. Пассивный MSAU соединяет порты внутренними связями так, чтобы станции образовали кольцо, и обеспечивает обход порта при выключении подключённого компьютера. Активный концентратор выполняет регенерацию сигналов и ресинхронизацию. Используются кабели STP Type 1, UTP категорий 3, 4, 5, а также волоконно-оптический кабель. Максимальное количество станций в кольце — 260, максимальная длина кольца — 4000 м. При использовании STP Type 1 длина ответвительных кабелей до 100 м, при UTP — до 45 м. Расстояние между пассивными концентраторами MSAU — до 100 м (STP) или 45 м (UTP); между активными — до 730 м (STP) или 365 м (UTP).
\textit{(Источник: Олифер, гл. 7.)}

\subsubsection{Технология FDDI}

\paragraph{Общая характеристика.}
FDDI (Fiber Distributed Data Interface) — технология локальных сетей со скоростью 100 Мбит/с, разработанная ANSI. Она использует двойное волоконно-оптическое кольцо, маркерный метод доступа и обеспечивает высокую отказоустойчивость. Максимальная длина кольца — около 100 км, число станций двойного подключения — до 500. FDDI также поддерживает вариант TP-PMD для витой пары категории 5. Технология FDDI во многом основывается на Token Ring, развивая и совершенствуя её идеи: кольцевая топология, маркерный доступ, приоритеты, но с более высокой скоростью и отказоустойчивостью.
\textit{(Источник: Олифер, гл. 7.)}

\paragraph{Двойное кольцо и отказоустойчивость.}
Сеть FDDI строится на основе двух оптоволоконных колец: первичного (Primary) и вторичного (Secondary). В нормальном режиме данные передаются по первичному кольцу; вторичное кольцо не используется. При обрыве кабеля или отказе станции выполняется операция «свёртывания» колец (wrap): первичное и вторичное кольца объединяются, образуя единое кольцо, и сеть продолжает работать. Станции могут иметь двойное подключение (Dual Attachment, DA) — к обоим кольцам, или одиночное (Single Attachment, SA) — только к первичному. Концентраторы также делятся на DAC (Dual Attachment Concentrator) и SAC (Single Attachment Concentrator). Станции DAS оснащаются оптическими обходными переключателями (Optical Bypass Switch), которые создают обходной путь при потере питания. Поддерживается подключение Dual Homing — двойное подключение к двум портам M одного или двух концентраторов, что повышает надёжность.
\textit{(Источник: Олифер, гл. 7.)}

\paragraph{Метод доступа FDDI.}
FDDI использует маркерный метод доступа. Для синхронного трафика станция всегда имеет право захватить маркер при его поступлении на фиксированное время. Для асинхронного трафика станция может захватить маркер только в том случае, если маркер выполнил оборот по кольцу достаточно быстро, что говорит об отсутствии перегрузок кольца. Станция измеряет время оборота маркера (Token Rotation Time, TRT) и сравнивает его с максимально допустимым временем оборота \(T_{\text{opr}}\). Если \(TRT < T_{\text{opr}}\), станция может захватить маркер и передать асинхронные кадры в течение времени \(T_{\text{opr}} - TRT\). Если кольцо перегружено и маркер опоздал, станция не имеет права захватывать маркер для асинхронного трафика. Такой метод отдаёт предпочтение синхронному трафику и регулирует загрузку кольца. Размер окна и время удержания маркера динамически адаптируются. Механизм приоритетов кадров, аналогичный Token Ring, в FDDI отсутствует; трафик делится на два класса — синхронный и асинхронный.
\textit{(Источник: Олифер, гл. 7.)}

\paragraph{Формат кадра FDDI.}
Формат кадра FDDI близок к формату кадра Token Ring. Кадр содержит начальный ограничитель, поле управления доступом, адреса назначения и источника, поле данных, контрольную сумму, конечный ограничитель и поле статуса. Основные отличия заключаются в отсутствии полей приоритетов. Признаки распознавания адреса, копирования кадра и ошибки позволяют сохранить процедуры обработки кадров станцией-отправителем, промежуточными станциями и станцией-получателем. Управление станцией выполняет уровень SMT (Station Management), который отвечает за управление и мониторинг всех остальных уровней стека FDDI. Узлы FDDI обмениваются специальными кадрами SMT для управления сетью.
\textit{(Источник: Олифер, гл. 7.)}

\paragraph{Физический уровень FDDI.}
Физический уровень FDDI разделён на два подуровня: независимый от среды PHY и зависящий от среды PMD. Технология FDDI поддерживает два подуровня PMD: для волоконно-оптического кабеля и для неэкранированной витой пары категории 5 (TP-PMD). В качестве основной физической среды используется многомодовый волоконно-оптический кабель 62,5/125 мкм. Для передачи света применяется длина волны 1300 нм. Предельное расстояние между узлами для многомодового кабеля — 2 км, для одномодового — до 10–40 км в зависимости от качества кабеля. Для витой пары TP-PMD максимальное расстояние — 100 м. Используется логическое кодирование 4B/5B в сочетании с физическим кодированием NRZI; тактовая частота — 125 МГц. Служебные символы кода 4B/5B, в частности Idle, постоянно передаются между портами в паузах, что позволяет контролировать состояние физических соединений.
\textit{(Источник: Олифер, гл. 7.)}

\subsubsection{Сравнение архитектур Ethernet, Token Ring и FDDI}

\begin{table}[h]
\centering
\caption{Сравнительная характеристика базовых технологий локальных сетей}
\label{tab:lan-arch-compare}
\begin{tabular}{p{2.5cm}p{3.5cm}p{3.5cm}p{3.5cm}}
\toprule
\textbf{Параметр} & \textbf{Ethernet} & \textbf{Token Ring} & \textbf{FDDI} \\
\midrule
Метод доступа & CSMA/CD (случайный) & Маркерный & Маркерный \\
Скорость & 10, 100, 1000, 10000 Мбит/с & 4, 16 Мбит/с & 100 Мбит/с \\
Логическая топология & Шина, звезда, дерево & Кольцо & Двойное кольцо \\
Физическая топология & Шина, звезда, дерево & Звезда (MSAU) & Звезда, кольцо \\
Среда передачи & Коаксиал, витая пара, оптоволокно & Витая пара, оптоволокно & Оптоволокно, витая пара (TP-PMD) \\
Отказоустойчивость & Зависит от топологии и коммутации & Активный монитор, beaconing & Свёртывание двойного кольца \\
Максимальное число станций & Ограничено доменом коллизий & 260 & 500 (DAS) \\
Максимальная длина & Ограничена диаметром домена коллизий & 4000 м & 100 км \\
Приоритеты & Нет (в базовом стандарте) & Да (0–7) & Нет (синхронный/асинхронный трафик) \\
\bottomrule
\end{tabular}
\end{table}

\subsubsection{Краткое резюме подраздела}
Ethernet использует случайный метод доступа CSMA/CD, MAC-адресацию и семейство скоростей от 10 Мбит/с до 10 Гбит/с; в коммутируемых сетях широко применяется полнодуплексный режим. Token Ring использует маркерный доступ, логическое кольцо, активный монитор, приоритеты и процедуры диагностики. FDDI обеспечивает 100 Мбит/с, двойное кольцо, маркерный доступ, высокую отказоустойчивость и поддержку как синхронного, так и асинхронного трафика. Архитектуры этих технологий отражают разные подходы к разделению среды, детерминированности доступа, надёжности и масштабируемости локальных сетей.
\textit{(Обобщение по: Олифер, гл. 7–8.)}

\subsection{Сеть Internet, доменная организация, семейство протоколов TCP/IP}

\noindent\textit{Основные источники: Олифер В.Г., Олифер Н.А. Компьютерные сети. Принципы, технологии, протоколы. 2-е изд. Гл. 11–14.}

\subsubsection{Internet как глобальная составная сеть}

\paragraph{Понятие составной сети и межсетевого взаимодействия.}
Internet представляет собой глобальную составную сеть (internetwork), объединяющую множество разнородных сетей — локальных, городских, ведомственных, операторских. Отдельные сети, входящие в Internet, называются подсетями (subnet). Они соединяются между собой маршрутизаторами, которые обеспечивают передачу пакетов между сетями с различными технологиями канального уровня. Организация совместной транспортной службы в такой неоднородной среде называется межсетевым взаимодействием (internetworking). Для универсальной адресации и маршрутизации в Internet используется стек протоколов TCP/IP, который обеспечивает независимость от конкретных технологий подсетей.
\textit{(Источник: Олифер, гл. 11.)}

\paragraph{Автономные системы.}
Internet состоит не только из сетей, но и из более крупных объединений — автономных систем (Autonomous System, AS). Автономная система — это совокупность сетей под единым административным управлением, обеспечивающим общую для всех входящих в неё маршрутизаторов политику маршрутизации. Обычно автономной системой управляет один поставщик услуг Интернета или крупная организация. Автономные системы соединяются внешними маршрутизаторами (exterior gateway). Между автономными системами используется единый внешний протокол маршрутизации — в настоящее время BGP-4. Внутри автономной системы могут применяться различные внутренние протоколы маршрутизации, например RIP или OSPF. Деление на автономные системы обеспечивает многоуровневый подход к маршрутизации: сначала выбирается маршрут на уровне автономных систем, затем — на уровне сетей внутри автономной системы.
\textit{(Источник: Олифер, гл. 14.)}

\paragraph{Структура Internet.}
Современный Internet представляет собой совокупность взаимосвязанных автономных систем с произвольной топологией связей. Магистральные операторы связи образуют ядро сети, к которому подключаются региональные и локальные операторы, корпоративные сети и сети доступа. На границе сети оператора устанавливаются пограничные маршрутизаторы, поддерживающие BGP. Внутри сети оператора работают внутренние протоколы маршрутизации. Такая иерархическая структура обеспечивает масштабируемость, гибкость и возможность подключения новых сетей без перестройки всей Internet.
\textit{(Источник: Олифер, гл. 14.)}

\subsubsection{Доменная организация и система доменных имён}

\paragraph{Символьные имена и проблема уникальности.}
Для пользователей IP-адреса неудобны, поэтому в сетях TCP/IP применяются символьные доменные имена. Доменная система имён имеет иерархическую древовидную структуру, допускающую использование в имени произвольного количества составных частей. Дерево имён начинается с корня, обозначаемого точкой. Старшая часть имени (домен верхнего уровня) определяет страну или тип организации; далее следуют поддомены и, наконец, имя конкретного хоста. Составные части доменного имени отделяются друг от друга точкой. Пример: \texttt{www1.zil.mmt.ru}. Разделение имени на части позволяет разделить административную ответственность за назначение уникальных имён между различными организациями. Совокупность имён, у которых несколько старших составных частей совпадают, образует домен. Домен, входящий в другой домен, называется поддоменом.
\textit{(Источник: Олифер, гл. 12.)}

\paragraph{Домены верхнего уровня.}
Домены верхнего уровня назначаются для каждой страны, а также на организационной основе. Для обозначения стран используются двухбуквенные аббревиатуры: \texttt{ru} (Россия), \texttt{uk} (Великобритания), \texttt{us} (Соединённые Штаты) и др. Для различных типов организаций используются следующие обозначения: \texttt{com} — коммерческие организации, \texttt{edu} — образовательные, \texttt{gov} — правительственные, \texttt{org} — некоммерческие, \texttt{net} — организации поддержки сетей. Каждый домен администрируется отдельной организацией, которая может разбивать его на поддомены и передавать функции администрирования другим организациям. Регистрация доменных имён в зоне \texttt{ru} осуществляется через аккредитованных регистраторов, координируемых РосНИИРОС.
\textit{(Источник: Олифер, гл. 12.)}

\paragraph{Система доменных имён DNS.}
Соответствие между доменными именами и IP-адресами устанавливается с помощью системы доменных имён (Domain Name System, DNS). DNS — это централизованная служба, основанная на распределённой базе отображений «доменное имя — IP-адрес». Служба DNS использует модель «клиент-сервер». DNS-серверы поддерживают распределённую базу отображений, а DNS-клиенты обращаются к серверам с запросами о разрешении доменного имени в IP-адрес. Для каждого домена имён создаётся свой DNS-сервер. Сервер домена хранит отображения для имён, которые заканчиваются на следующем ниже уровне иерархии, а также ссылки на DNS-серверы своих поддоменов. Ссылки представляют собой IP-адреса соответствующих серверов. Для обслуживания корневого домена выделено несколько дублирующих друг друга DNS-серверов.
\textit{(Источник: Олифер, гл. 12.)}

\paragraph{Разрешение имён.}
Существует две основные схемы разрешения DNS-имён. При нерекурсивной (итеративной) схеме DNS-клиент сам выполняет последовательность запросов к разным серверам: сначала к корневому, затем к серверу домена верхнего уровня, затем к серверу поддомена и т. д., пока не найдёт сервер, хранящий соответствие запрошенного имени IP-адресу. При рекурсивной схеме клиент запрашивает локальный DNS-сервер, который при необходимости выполняет итеративные запросы к другим серверам и возвращает клиенту готовый ответ. Практически все DNS-клиенты используют рекурсивную процедуру. Для ускорения поиска DNS-серверы широко применяют кэширование проходящих через них ответов. Время кэширования обычно составляет от нескольких часов до нескольких дней.
\textit{(Источник: Олифер, гл. 12.)}

\subsubsection{Семейство протоколов TCP/IP}

\paragraph{Общая характеристика стека TCP/IP.}
Стек TCP/IP был разработан по инициативе Министерства обороны США для связи экспериментальной сети ARPAnet с другими сетями как набор общих протоколов для разнородной вычислительной среды. Сегодня стек TCP/IP используется для связи компьютеров всемирной информационной сети Internet, а также в огромном числе корпоративных сетей. Стек TCP/IP имеет четыре уровня: прикладной, транспортный, уровень межсетевого взаимодействия и уровень сетевых интерфейсов. Такое деление не полностью совпадает с семиуровневой моделью OSI, но функционально покрывает все необходимые задачи межсетевого взаимодействия.
\textit{(Источник: Олифер, гл. 11.)}

\begin{table}[h]
\centering
\caption{Уровни стека TCP/IP и их основные протоколы}
\label{tab:tcpip-levels-detail}
\begin{tabular}{p{4cm}p{9cm}}
\toprule
\textbf{Уровень TCP/IP} & \textbf{Основные протоколы и функции} \\
\midrule
Прикладной & HTTP, FTP, SMTP, DNS, DHCP, SNMP, telnet; сетевые службы и приложения \\
Транспортный & TCP — надёжная доставка с установлением соединения; UDP — дейтаграммная доставка \\
Межсетевого взаимодействия & IP, ICMP, ARP; маршрутизация RIP, OSPF, BGP \\
Сетевых интерфейсов & Ethernet, Token Ring, FDDI, PPP, SLIP, Frame Relay, ATM \\
\bottomrule
\end{tabular}
\end{table}

\paragraph{Прикладной уровень.}
Прикладной уровень стека TCP/IP объединяет службы, предоставляемые системой пользовательским приложениям. К прикладным протоколам относятся: протокол копирования файлов FTP, протокол эмуляции терминала telnet, простой протокол передачи электронной почты SMTP, протокол передачи гипертекстовой информации HTTP, система доменных имён DNS, протокол динамического конфигурирования узлов DHCP, протокол управления сетями SNMP. Протоколы прикладного уровня устанавливаются на хостах и реализуются программными системами, построенными в архитектуре «клиент-сервер». Они отрабатывают логику приложений и обращаются к протоколам нижних уровней как к набору инструментов.
\textit{(Источник: Олифер, гл. 11.)}

\paragraph{Транспортный уровень.}
Транспортный уровень стека TCP/IP может предоставлять вышележащему уровню два типа сервиса: гарантированную доставку (протокол TCP) и доставку «по возможности» (протокол UDP). Протокол TCP обеспечивает надёжную доставку данных, устанавливая логическое соединение, нумеруя сегменты, подтверждая их приём квитанциями, организуя повторные передачи при потере, распознавая и уничтожая дубликаты и доставляя данные в том порядке, в котором они были отправлены. UDP является простейшим дейтаграммным протоколом, который используется в том случае, когда задача надёжного обмена данными либо вообще не ставится, либо решается средствами более высокого уровня. Протоколы TCP и UDP устанавливаются на хостах.
\textit{(Источник: Олифер, гл. 11, 13.)}

\paragraph{Уровень межсетевого взаимодействия.}
Уровень межсетевого взаимодействия (internet) обеспечивает перемещение пакетов в пределах всей составной сети. Основным протоколом этого уровня является протокол IP. В его задачу входит продвижение пакета между подсетями — от одного пограничного маршрутизатора до другого, пока пакет не попадёт в сеть назначения. Протокол IP устанавливается не только на хостах, но и на всех маршрутизаторах. IP — дейтаграммный протокол, работающий без установления соединений по принципу «по возможности»; он не берёт на себя ответственность за доставку пакета до узла назначения. На этом же уровне работают протокол межсетевых управляющих сообщений ICMP, протокол разрешения адресов ARP, а также протоколы маршрутизации RIP, OSPF, BGP.
\textit{(Источник: Олифер, гл. 11–14.)}

\paragraph{Уровень сетевых интерфейсов.}
Уровень сетевых интерфейсов отвечает за организацию интерфейса с технологиями подсетей. Он обеспечивает инкапсуляцию IP-пакетов в кадры технологий канального уровня: Ethernet, Token Ring, FDDI, PPP, SLIP, Frame Relay, ATM и др. Перемещение пакета IP можно рассматривать как последовательность «прыжков» от одного маршрутизатора до другого. На каждом шаге протоколы TCP/IP обращаются к транспортным средствам соответствующей подсети. Уровень сетевых интерфейсов не регламентируется, но поддерживает все популярные стандарты физического и канального уровней.
\textit{(Источник: Олифер, гл. 11.)}

\subsubsection{Протокол IP}

\paragraph{Функции и особенности.}
Протокол IP относится к протоколам без установления соединения. Он обрабатывает каждый IP-пакет как независимую единицу, не имеющую связи ни с какими другими IP-пакетами. В протоколе IP нет механизмов квитирования, упорядочивания, повторных передач или других подобных функций. Если во время продвижения пакета произошла ошибка, протокол IP по своей инициативе ничего не предпринимает для её исправления. Все вопросы обеспечения надёжности при доставке данных решает протокол TCP. Важной особенностью протокола IP является его способность выполнять динамическую фрагментацию пакетов при передаче их между сетями с различными значениями MTU.
\textit{(Источник: Олифер, гл. 13.)}

\paragraph{Структура IP-пакета.}
IP-пакет состоит из заголовка и поля данных. Заголовок обычно имеет длину 20 байт и содержит следующие поля: номер версии (Version), длина заголовка (IHL), тип сервиса (Type of Service, ToS), общая длина (Total Length), идентификатор пакета (Identification), флаги (Flags), смещение фрагмента (Fragment Offset), время жизни (Time to Live), протокол верхнего уровня (Protocol), контрольная сумма заголовка (Header Checksum), IP-адрес источника (Source IP Address), IP-адрес назначения (Destination IP Address), параметры (IP Options), выравнивание (Padding). Поле «Протокол верхнего уровня» указывает, какому протоколу принадлежит информация в поле данных: 6 — TCP, 17 — UDP, 1 — ICMP. Максимальная длина IP-пакета — 65 535 байт.
\textit{(Источник: Олифер, гл. 13.)}

\paragraph{IP-адресация.}
IP-адрес имеет длину 4 байта и состоит из номера сети и номера узла. Для определения границы между номером сети и номером узла используются два подхода: классы адресов и маски. Класс адреса определяется значениями нескольких первых битов. В адресах класса A под номер сети отводится один байт, а остальные три — под номер узла; в адресах класса B — по два байта; в адресах класса C — три байта под номер сети и один байт под номер узла. Классы D и E используются для групповой адресации и зарезервированы. При использовании масок граница между номером сети и номером узла задаётся маской — числом, двоичная запись которого содержит единицы в тех разрядах, которые интерпретируются как номер сети. Маски позволяют гибко структурировать адресное пространство, разбивать сеть на подсети и агрегировать маршруты (технология CIDR).
\textit{(Источник: Олифер, гл. 12.)}

\paragraph{Фрагментация.}
Фрагментация — это деление поля данных исходного пакета на части и оформление этих частей в виде пакетов меньшего размера — фрагментов. Фрагментация выполняется при невозможности передать пакет в следующую по маршруту сеть из-за того, что его размер превосходит максимально допустимый размер единицы передаваемых данных в этой сети. Все фрагменты перемещаются по составной сети независимо друг от друга. Сбор исходного пакета из фрагментов осуществляется модулем IP в узле назначения (и никогда на промежуточных маршрутизаторах). При сборке устанавливается таймер, ограничивающий время ожидания фрагментов. Если ко времени истечения таймера хотя бы один фрагмент не дойдёт до узла назначения, все фрагменты отбрасываются.
\textit{(Источник: Олифер, гл. 13.)}

\subsubsection{Протоколы TCP и UDP}

\paragraph{Протокол TCP.}
Протокол TCP работает непосредственно над протоколом IP и использует для транспортировки своих блоков данных этот дейтаграммный и потенциально ненадёжный протокол. Надёжность передачи данных протоколом TCP достигается за счёт установления логических соединений между взаимодействующими процессами. Информация, поступающая к протоколу TCP в рамках логического соединения от протоколов более высокого уровня, рассматривается как неструктурированный поток байтов. Поступающие данные буферизуются; для передачи на сетевой уровень из буфера «вырезается» некоторая непрерывная часть данных — сегмент. В заголовок сегмента TCP помещает номер порта назначения и номер порта отправителя. При установлении соединения стороны договариваются о максимальном размере сегмента, максимальном объёме данных, которые можно передать без подтверждения, и начальном порядковом номере байта. В рамках соединения правильность передачи каждого сегмента подтверждается квитанцией получателя. Используется алгоритм скользящего окна, при этом окно определено на множестве нумерованных байтов потока данных.
\textit{(Источник: Олифер, гл. 13.)}

\paragraph{Протокол UDP.}
UDP является простейшим дейтаграммным протоколом. Он используется в том случае, когда задача надёжного обмена данными либо вообще не ставится, либо решается средствами более высокого уровня — системными прикладными службами или пользовательскими приложениями. UDP не устанавливает соединения, не подтверждает приём данных и не выполняет повторных передач. Заголовок UDP содержит номера портов отправителя и получателя, длину и контрольную сумму. UDP обеспечивает минимальные накладные расходы и используется для служб, где важна простота и низкая задержка, например для DNS, DHCP, SNMP, потокового мультимедиа.
\textit{(Источник: Олифер, гл. 13.)}

\paragraph{Порты и сокеты.}
Каждый компьютер может выполнять несколько процессов, более того, прикладной процесс может иметь несколько точек входа, выступающих в качестве адресов назначения для пакетов данных. Пакеты, поступающие на транспортный уровень, организуются операционной системой в виде множества очередей к точкам входа различных прикладных процессов. В терминологии TCP/IP такие системные очереди называются портами. Порт однозначно определяет приложение в пределах компьютера. Номер порта в совокупности с номером сети и номером конечного узла однозначно определяют прикладной процесс в сети. Этот набор идентифицирующих параметров (IP-адрес, номер порта) имеет название сокет (socket). Каждый взаимодействующий процесс идентифицируется сокетом, а каждое соединение — парой сокетов взаимодействующих процессов.
\textit{(Источник: Олифер, гл. 13.)}

\subsubsection{Служебные протоколы и маршрутизация}

\paragraph{Протокол ARP.}
Для определения локального адреса по IP-адресу используется протокол разрешения адресов ARP. В сетях с широковещанием (Ethernet, Token Ring, FDDI) ARP-запрос рассылается широковещательно; интерфейс, опознавший свой IP-адрес, формирует ARP-ответ, указывая свой IP-адрес и локальный адрес. Для ускорения процедуры преобразования адресов протокол ARP кэширует полученные ответы в ARP-таблицах. В сетях, не поддерживающих широковещание, ARP-таблицы могут храниться централизованно на выделенном ARP-сервере или формироваться вручную.
\textit{(Источник: Олифер, гл. 12.)}

\paragraph{Протокол ICMP.}
Протокол межсетевых управляющих сообщений ICMP предназначен для обмена информацией об ошибках между маршрутизаторами сети и узлом-источником пакета. С помощью специальных пакетов ICMP сообщает о невозможности доставки пакета, о превышении времени жизни или продолжительности сборки пакета из фрагментов, об аномальных величинах параметров, об изменении маршрута пересылки и типа обслуживания, о состоянии системы. ICMP также используется для диагностики сети (утилиты ping, traceroute).
\textit{(Источник: Олифер, гл. 13.)}

\paragraph{Протоколы маршрутизации.}
Для автоматического построения таблиц маршрутизации маршрутизаторы обмениваются информацией о топологии составной сети с помощью протоколов маршрутизации. Протокол RIP является дистанционно-векторным протоколом и использует метрику в хопах. Протокол OSPF основан на алгоритме состояния связей, поддерживает области, метрики пропускной способности и балансировку загрузки. Протокол BGP применяется между автономными системами и отвечает за выбор маршрута как последовательности автономных систем. Внутри автономной системы могут использоваться RIP, OSPF, IS-IS и другие протоколы. Маршрутизаторы хранят таблицы маршрутизации, в которых указываются сети назначения, адреса следующих маршрутизаторов и метрики.
\textit{(Источник: Олифер, гл. 14.)}

\subsubsection{Прикладные протоколы}

\paragraph{Основные прикладные протоколы.}
К прикладным протоколам стека TCP/IP относятся: FTP (передача файлов), telnet (эмуляция терминала), SMTP (передача электронной почты), HTTP (передача гипертекста), DNS (разрешение доменных имён), DHCP (автоматическое назначение IP-адресов), SNMP (управление сетевыми устройствами). Протокол FTP предназначен для обмена файлами между компьютерами сети. Протокол telnet позволяет пользователю удалённо работать с компьютером так, как если бы он управлял им с помощью локального терминала. Протокол SMTP используется для передачи электронной почты. Протокол HTTP является основным протоколом службы WWW. DNS устанавливает соответствие между доменными именами и IP-адресами. DHCP автоматизирует процесс назначения IP-адресов. SNMP используется для управления сетевыми устройствами и получения от них информации об их статусе и производительности.
\textit{(Источник: Олифер, гл. 11–12, 20.)}

\subsubsection{Краткое резюме подраздела}
Internet представляет собой глобальную составную сеть, объединяющую автономные системы и отдельные сети. Доменная система имён DNS обеспечивает иерархическое символьное именование и разрешение имён в IP-адреса. Стек TCP/IP включает прикладной, транспортный, межсетевой уровни и уровень сетевых интерфейсов. Ключевыми протоколами являются IP, TCP, UDP, ICMP, ARP, DNS, DHCP и протоколы маршрутизации RIP, OSPF, BGP. Прикладные протоколы HTTP, FTP, SMTP, telnet, SNMP обеспечивают работу сетевых служб. Совместное применение этих протоколов и механизмов обеспечивает масштабируемость, надёжность и совместимость глобальной сети Internet.
\textit{(Обобщение по: Олифер, гл. 11–14.)}

\subsection{Информационно-вычислительные сети и распределенная обработка информации}

\noindent\textit{Основные источники: Олифер В.Г., Олифер Н.А. Компьютерные сети. Принципы, технологии, протоколы. 2-е изд. Гл. 1–2, 15.}

\subsubsection{Понятие распределённой вычислительной системы}

\paragraph{Определение распределённой системы.}
Компьютерные сети представляют собой частный случай распределённых вычислительных систем, в которых группа компьютеров согласованно выполняет набор взаимосвязанных задач, обмениваясь данными в автоматическом режиме. Наряду с компьютерными сетями к распределённым системам относят также мультипроцессорные компьютеры и многомашинные вычислительные комплексы (кластеры). Общим признаком распределённой системы является наличие нескольких центров обработки данных, которые взаимодействуют между собой для решения общей задачи. Степень связанности этих центров и способ организации взаимодействия определяют тип распределённой системы.
\textit{(Источник: Олифер, гл. 2.)}

\paragraph{Мультипроцессорные компьютеры.}
В мультипроцессорных компьютерах имеется несколько процессоров, каждый из которых может относительно независимо от остальных выполнять свою программу. Все периферийные устройства являются для всех процессоров мультипроцессорной системы общими. В мультипроцессоре существует общая для всех процессоров операционная система, которая оперативно распределяет вычислительную нагрузку между процессорами. Взаимодействие между отдельными процессорами организуется наиболее простым способом — через общую оперативную память. Мультипроцессору не свойственна территориальная распределённость — все его блоки располагаются в одном или нескольких близко расположенных конструктивах, как и у обычного компьютера. Основное достоинство мультипроцессора — высокая производительность, достигаемая за счёт параллельной работы нескольких процессоров. При наличии общей памяти взаимодействие процессоров происходит очень быстро, поэтому мультипроцессоры могут эффективно выполнять даже приложения с высокой степенью связи по данным. Ещё одним важным свойством мультипроцессорных систем является отказоустойчивость — способность продолжать работу при отказах некоторых элементов, например процессоров или блоков памяти. При этом производительность снижается, но не до нуля.
\textit{(Источник: Олифер, гл. 2.)}

\paragraph{Кластеры (многомашинные вычислительные комплексы).}
Кластер — это вычислительный комплекс, состоящий из нескольких компьютеров, каждый из которых работает под управлением собственной операционной системы, а также программных и аппаратных средств связи компьютеров, которые обеспечивают работу всех компьютеров комплекса как единого целого. В отличие от мультипроцессора, в котором избыточность реализована на уровне процессорных блоков, кластер состоит из нескольких законченных, способных работать автономно, как правило, стандартных компьютеров, каждый из которых имеет обычную структуру, включающую один или несколько процессорных блоков, оперативную память и периферийные устройства. Благодаря специальному программному и аппаратному обеспечению межкомпьютерных связей для пользователя кластер выглядит как единый компьютер. При этом каждый компьютер, называемый также узлом кластера, может быть как однопроцессорным, так и мультипроцессорным — на организацию кластера это влияния не оказывает. Кластеры применяют для повышения надёжности и производительности вычислительной системы. Надёжность повышается за счёт того, что при отказе одного из узлов кластера вычислительная нагрузка переносится на другой узел. Для выполнения этой операции в кластере используется два типа связей между узлами: межпроцессорные связи и связи за счёт разделяемых дисков. Межпроцессорные связи используются узлами для обмена служебной информацией. С помощью этих связей каждый узел кластера периодически проверяет состояние других узлов и выполняемых ими вычислительных задач. Если какой-либо узел или одна из его задач изменили своё состояние на неработоспособное, то начинается процедура перемещения нагрузки на один из работоспособных узлов. В этой процедуре важную роль играют разделяемые диски. Защищаемая задача должна хранить свои данные на одном из таких дисков, чтобы новый узел смог продолжить их использовать после отказа основного. Так как надёжность дисковых накопителей достаточно высока, то существенно повышается и надёжность кластера по сравнению с отдельным компьютером. Время перевода нагрузки на другой узел кластера при отказе значительно больше, чем время перехода на другой процессор в мультипроцессоре, так как оно связано с активизацией новой копии программного процесса на другом узле. При этом также возможна потеря части данных, находившихся в оперативной памяти отказавшего узла, но для определённых типов вычислительной нагрузки, например систем управления базами данных или web-серверов, эти потери не сказываются на возможности продолжать вычисления. Если кластер применяется для повышения производительности, то каждая задача распараллеливается на несколько ветвей, которые выполняются одновременно на нескольких узлах кластера. Синхронизация работы нескольких копий задачи или их ветвей, а также синхронизация вырабатываемых ими данных осуществляются как за счёт межпроцессорных связей, так и за счёт разделяемой дисковой памяти. Менее тесные и менее скоростные связи между узлами кластера по сравнению со связями процессоров в мультипроцессоре диктуют область применения кластеров — это задачи, достаточно независимые по данным. Для организации межпроцессорных связей в кластерах часто используются специализированные технологии, приспособленные к решению специфических задач взаимодействия компьютеров в кластере. Однако в последнее время всё чаще для этой цели применяются стандартные технологии локальных сетей, например Fast Ethernet и Gigabit Ethernet.
\textit{(Источник: Олифер, гл. 2.)}

\subsubsection{Компьютерная сеть как распределённая система}

\paragraph{Особенности сетей.}
В вычислительных сетях программные и аппаратные связи являются ещё более слабыми, а автономность обрабатывающих блоков проявляется в наибольшей степени. Основными элементами сети являются стандартные компьютеры, не имеющие ни общих блоков памяти, ни совместно управляемых дисковых накопителей. Связь между компьютерами осуществляется с помощью специальных периферийных устройств — сетевых адаптеров, соединённых относительно протяжёнными каналами связи. Каждый компьютер работает под управлением собственной операционной системы, а какая-либо «общая» операционная система, распределяющая работу между компьютерами сети, отсутствует. Взаимодействие между компьютерами сети происходит путём передачи сообщений через сетевые адаптеры и каналы связи. С помощью этих сообщений один компьютер обычно запрашивает доступ к локальным ресурсам другого компьютера. Таким ресурсом может являться, например, вычислительная мощность компьютера в целом, именно этот ресурс использует удалённый пользователь, когда запускает свою программу на «чужом» компьютере. В качестве совместно используемых ресурсов часто выступают данные, хранящиеся на диске, а также разнообразные периферийные устройства — принтеры, модемы, факс-аппараты и т. д.
\textit{(Источник: Олифер, гл. 2.)}

\paragraph{Разделение ресурсов и клиент-серверная модель.}
Разделение локальных ресурсов каждого компьютера между всеми пользователями сети — основная цель создания вычислительной сети. Для этого недостаточно снабдить компьютеры сетевыми адаптерами и соединить их кабельной системой, необходимы ещё некоторые добавления к их операционным системам. На тех компьютерах, ресурсы которых должны быть доступны всем пользователям сети, необходимо добавить модули, которые постоянно будут находиться в режиме ожидания запросов, поступающих по сети от других компьютеров. Обычно такие модули называются программными серверами (server), так как их главная задача — обслуживать (serve) запросы на доступ к ресурсам. На компьютерах, которые должны иметь доступ к ресурсам других компьютеров, необходимо добавить программные модули, которые формируют запросы на доступ к удалённым ресурсам и передают их в сеть. Такие модули называются клиентами (client). Набор модулей «клиент — сервер» представляет собой распределённую программу, реализующую сетевую службу (service). Термины «клиент» и «сервер» используются не только для обозначения программных модулей, но и компьютеров, подключённых к сети. Если компьютер предоставляет свои ресурсы другим компьютерам сети, то он называется сервером, а если он их потребляет — клиентом. Иногда один и тот же компьютер может одновременно играть роли и сервера, и клиента.
\textit{(Источник: Олифер, гл. 2.)}

\paragraph{Сетевые службы.}
Для конечного пользователя сеть — это прежде всего тот набор сетевых служб, с помощью которых он получает возможность просмотреть список имеющихся в сети компьютеров, прочитать удалённый файл, распечатать документ на «чужом» принтере или послать почтовое сообщение. Именно совокупность предоставляемых возможностей — насколько широк их выбор, насколько они удобны, надёжны и безопасны — определяет для пользователя облик той или иной сети. Кроме собственно обмена данными, сетевые службы должны решать и другие, более специфические задачи, например задачи, порождаемые распределённой обработкой данных. К таким задачам относится обеспечение непротиворечивости нескольких копий данных, размещённых на разных машинах (служба репликации), или организация выполнения одной задачи параллельно на нескольких машинах сети (служба вызова удалённых процедур). Среди сетевых служб можно выделить административные, которые в основном ориентированы не на простого пользователя, а на администратора и служат для организации правильной работы сети в целом. Служба администрирования учётных записей о пользователях, система мониторинга сети, служба безопасности — всё это примеры административных служб. Реализация сетевых служб осуществляется программными средствами. Основные службы — файловая служба и служба печати — обычно предоставляются сетевой операционной системой, а вспомогательные, например служба баз данных, факса или передачи голоса, — системными сетевыми приложениями, или утилитами, работающими в тесном контакте с сетевой ОС. Вообще говоря, распределение служб между ОС и утилитами достаточно условно и меняется в конкретных реализациях ОС.
\textit{(Источник: Олифер, гл. 2.)}

\subsubsection{Распределённые приложения}

\paragraph{Определение и отличие от локальных приложений.}
Распределённое приложение состоит из нескольких частей, каждая из которых выполняет какую-то определённую законченную работу по решению прикладной задачи. Например, одна часть приложения, выполняющаяся на компьютере пользователя, может поддерживать специализированный графический интерфейс, вторая — работать на мощном выделенном компьютере и заниматься статистической обработкой введённых пользователем данных, третья — заносить полученные результаты в базу данных на компьютере с установленной стандартной СУБД. Распределённые приложения в полной мере используют потенциальные возможности распределённой обработки, предоставляемые вычислительной сетью, и поэтому часто называются сетевыми приложениями. Следует подчеркнуть, что не всякое приложение, выполняемое в сети, является сетевым. Существует большое количество популярных приложений, которые не являются распределёнными и целиком выполняются на одном компьютере сети. Тем не менее и такие приложения могут использовать преимущества сети за счёт встроенных в операционную систему сетевых служб. Значительная часть истории локальных сетей связана как раз с использованием таких нераспределённых приложений. Рассмотрим, например, как происходила работа пользователя с известной в своё время СУБД dBase. Обычно файлы базы данных, с которыми работали все пользователи сети, располагались на файловом сервере. Сама же СУБД хранилась на каждом клиентском компьютере в виде единого программного модуля. Программа dBase была рассчитана на обработку только локальных данных, то есть данных, расположенных на том же компьютере, что и сама программа. Пользователь запускал dBase на своём компьютере, и программа искала данные на локальном диске, совершенно не принимая во внимание существование сети. Чтобы обрабатывать с помощью dBase данные, расположенные на удалённом компьютере, пользователь обращался к услугам файловой службы, которая доставляла данные с сервера на клиентский компьютер и создавала для СУБД эффект локального хранения. Большинство приложений, используемых в локальных сетях в середине 80-х годов, являлись обычными нераспределёнными приложениями. И это понятно — они были написаны для автономных компьютеров, а потом просто были перенесены в сетевую среду. Создание же распределённых приложений, хотя и сулило много преимуществ (снижение сетевого трафика, специализация компьютеров), оказалось делом совсем не простым. Нужно было решать множество дополнительных проблем — на сколько частей разбить приложение, какие функции возложить на каждую часть, как организовать взаимодействие этих частей, чтобы в случае сбоев и отказов оставшиеся части корректно завершали работу и т. д. и т. п. Поэтому до сих пор только небольшая часть приложений является распределёнными, хотя очевидно, что именно за этим классом приложений будущее, так как они в полной мере могут использовать потенциальные возможности сетей по распараллеливанию вычислений.
\textit{(Источник: Олифер, гл. 2.)}

\subsubsection{Репликация данных и удалённые вызовы}

\paragraph{Репликация.}
Одной из задач распределённой обработки является обеспечение непротиворечивости нескольких копий данных, размещённых на разных машинах. Эта задача решается службой репликации. Репликация повышает доступность данных и надёжность их хранения, но требует согласования копий при изменениях. Различают синхронную и асинхронную репликацию. При синхронной репликации все копии обновляются в рамках одной транзакции, что гарантирует их согласованность, но увеличивает задержки и требует постоянной связи между узлами. При асинхронной репликации обновления распространяются с задержкой, что повышает производительность и устойчивость к сбоям, но допускает временное расхождение копий. Выбор стратегии репликации зависит от требований приложения к согласованности, доступности и производительности.
\textit{(Источник: Олифер, гл. 2.)}

\paragraph{Удалённый вызов процедур.}
Для организации выполнения одной задачи параллельно на нескольких машинах сети используется служба вызова удалённых процедур (Remote Procedure Call, RPC). RPC позволяет программе вызывать процедуры, выполняющиеся на другом компьютере, так же, как если бы они были локальными. Механизм RPC скрывает от программиста детали сетевого взаимодействия: упаковку параметров, передачу запроса по сети, ожидание ответа и распаковку результата. RPC упрощает построение распределённых приложений, но требует решения проблем, связанных с надёжностью доставки, идентификацией удалённых процедур, обработкой ошибок и безопасностью.
\textit{(Источник: Олифер, гл. 2.)}

\subsubsection{Прозрачность и именование ресурсов}

\paragraph{Прозрачность доступа.}
При определении степени удобства разделяемого ресурса часто употребляют термин «прозрачность». Прозрачный доступ — это такой доступ, при котором пользователь не замечает, где расположен нужный ему ресурс — на собственном или удалённом компьютере. После того как пользователь смонтировал удалённую файловую систему в своё дерево каталогов, доступ к удалённым файлам становится для него совершенно прозрачным. Сама операция монтирования также может иметь разную степень прозрачности — в сетях с меньшей прозрачностью пользователь должен знать и задавать в команде имя компьютера, на котором хранится удалённая файловая система, в сетях с большей степенью прозрачности соответствующий программный компонент сети производит поиск разделяемых томов файлов безотносительно мест их хранения, а затем предоставляет их пользователю в удобном для него виде, например в виде списка или набора значков.
\textit{(Источник: Олифер, гл. 2.)}

\paragraph{Именование ресурсов.}
Для обеспечения прозрачности важен способ адресации (именования) разделяемых сетевых ресурсов. Имена разделяемых сетевых ресурсов не должны зависеть от их физического расположения на том или ином компьютере. В идеале пользователь не должен ничего менять в своей работе, если администратор сети переместил том или каталог с одного компьютера на другой. Сам администратор и сетевая операционная система имеют информацию о расположении файловых систем, но от пользователя она скрыта. Такая степень прозрачности пока редко встречается в сетях — обычно для получения доступа к ресурсам определённого компьютера сначала приходится устанавливать с ним логическое соединение. Такой подход применяется, например, в сетях Windows NT.
\textit{(Источник: Олифер, гл. 2.)}

\paragraph{Аспекты прозрачности.}
Концепция прозрачности может быть применена к различным аспектам сети. Прозрачность расположения означает, что от пользователя не требуется знаний о месте расположения программных и аппаратных ресурсов, таких как процессоры, принтеры, файлы и базы данных. Имя ресурса не должно включать информацию о месте его расположения, поэтому имена типа \texttt{mashine1:prog.c} или \texttt{\textbackslash\textbackslash ftp\_serv\textbackslash pub} прозрачными не являются. Аналогично, прозрачность перемещения означает, что ресурсы должны свободно перемещаться из одного компьютера в другой без изменения своих имён. Ещё одним из возможных аспектов прозрачности является прозрачность параллелизма, заключающаяся в том, что процесс распараллеливания вычислений происходит автоматически, без участия программиста, при этом система сама распределяет параллельные ветви приложения по процессорам и компьютерам сети. В настоящее время нельзя сказать, что свойство прозрачности в полной мере присуще многим вычислительным сетям, это скорее цель, к которой стремятся разработчики современных сетей.
\textit{(Источник: Олифер, гл. 2.)}

\subsubsection{Требования к ИВС и качество обслуживания}

\paragraph{Основные требования.}
Главным требованием, предъявляемым к сетям, является выполнение сетью того набора услуг, для оказания которых она предназначена: например, предоставление доступа к файловым архивам или страницам публичных web-сайтов, обмен электронной почтой в пределах предприятия или в глобальных масштабах, интерактивный обмен речевыми сообщениями IP-телефонии и т. п. Все остальные требования — производительность, надёжность, совместимость, управляемость, защищённость, расширяемость и масштабируемость — связаны с качеством выполнения основной задачи. И хотя все эти требования весьма важны, часто понятие «качество обслуживания» (Quality of Service, QoS) компьютерной сети трактуется более узко — в него включаются только две самые важные характеристики сети — производительность и надёжность.
\textit{(Источник: Олифер, гл. 2, 15.)}

\paragraph{Производительность.}
Потенциально высокая производительность — это одно из основных преимуществ распределённых систем, к которым относятся компьютерные сети. Это свойство обеспечивается принципиальной, но не всегда практически реализуемой возможностью распараллеливания работ между несколькими компьютерами сети. Существует несколько основных характеристик производительности сети: время реакции; скорость передачи данных; задержка передачи и вариация задержки передачи. Время реакции сети является интегральной характеристикой производительности сети с точки зрения пользователя. Именно эту характеристику имеет в виду пользователь, когда говорит: «Сегодня сеть работает медленно». В общем случае время реакции определяется как интервал времени между возникновением запроса пользователя к какой-либо сетевой службе и получением ответа на этот запрос. Очевидно, что значение этого показателя зависит от типа службы, к которой обращается пользователь, от того, какой пользователь и к какому серверу обращается, а также от текущего состояния элементов сети — загруженности сегментов, коммутаторов и маршрутизаторов, через которые проходит запрос, загруженности сервера и т. п. Поэтому имеет смысл использовать также и средневзвешенную оценку времени реакции сети, усредняя этот показатель по пользователям, серверам и времени дня. Скорость передачи данных отражает объём данных, переданных сетью или её частью в единицу времени. Пропускная способность уже не является пользовательской характеристикой, так как она говорит о скорости выполнения внутренних операций сети — передаче пакетов данных между узлами сети через различные коммуникационные устройства. Зато она непосредственно характеризует качество выполнения основной функции сети — транспортировки сообщений — и поэтому чаще используется при анализе производительности сети, чем время реакции. Скорость передачи данных измеряется либо в битах в секунду, либо в пакетах в секунду. Пропускная способность может быть мгновенной, максимальной и средней. Задержка передачи определяется как задержка между моментом поступления данных на вход какого-либо сетевого устройства или части сети и моментом появления их на выходе этого устройства. Этот параметр производительности по смыслу близок ко времени реакции сети, но отличается тем, что всегда характеризует только сетевые этапы обработки данных, без задержек обработки конечными узлами сети.
\textit{(Источник: Олифер, гл. 2.)}

\paragraph{Надёжность и безопасность.}
Одной из первоначальных целей создания распределённых систем, к которым относятся и вычислительные сети, являлось достижение большей надёжности по сравнению с отдельными вычислительными машинами. Важно различать несколько аспектов надёжности. Для технических устройств используются такие показатели надёжности, как среднее время наработки на отказ, вероятность отказа, интенсивность отказов. Однако эти показатели пригодны только для оценки надёжности простых элементов и устройств, которые могут находиться только в двух состояниях — работоспособном или неработоспособном. Сложные системы, состоящие из многих элементов, кроме состояний работоспособности и неработоспособности, могут иметь и другие промежуточные состояния, которые эти характеристики не учитывают. В связи с этим для оценки надёжности сложных систем применяется другой набор характеристик. Готовность, или коэффициент готовности (availability), означает долю времени, в течение которого система может быть использована. Готовность может быть улучшена путём введения избыточности в структуру системы: ключевые элементы системы должны существовать в нескольких экземплярах, чтобы при отказе одного из них функционирование системы обеспечивали другие. Чтобы компьютерную систему можно было отнести к высоконадёжным, она должна как минимум обладать высокой готовностью, но этого недостаточно. Необходимо обеспечить сохранность данных и защиту их от искажений. Кроме этого, должна поддерживаться согласованность (непротиворечивость) данных, например, если для повышения надёжности на нескольких файловых серверах хранится несколько копий данных, то нужно постоянно обеспечивать их идентичность. Так как сеть работает на основе механизма передачи пакетов между конечными узлами, то одной из характерных характеристик надёжности является вероятность доставки пакета узлу назначения без искажений. Наряду с этой характеристикой могут использоваться и другие показатели: вероятность потери пакета, вероятность искажения отдельного бита передаваемых данных, отношение потерянных пакетов к доставленным. Другим аспектом общей надёжности является безопасность (security), то есть способность системы защитить данные от несанкционированного доступа. В распределённой системе это сделать гораздо сложнее, чем в централизованной. В сетях сообщения передаются по линиям связи, часто проходящим через общедоступные помещения, в которых могут быть установлены средства прослушивания линий. Другим уязвимым местом могут быть оставленные без присмотра персональные компьютеры. Кроме того, всегда имеется потенциальная угроза взлома защиты сети от неавторизованных пользователей, если сеть имеет выходы в глобальные сети общего пользования. Ещё одной характеристикой надёжности является отказоустойчивость (fault tolerance). В сетях под отказоустойчивостью понимается способность системы скрывать от пользователя отказ отдельных её элементов. Например, если копии таблицы базы данных хранятся одновременно на нескольких файловых серверах, то пользователи могут просто не заметить отказ одного из них. В отказоустойчивой системе отказ одного из её элементов приводит к некоторому снижению качества её работы (деградации), а не к полному останову.
\textit{(Источник: Олифер, гл. 2.)}

\paragraph{Расширяемость и масштабируемость.}
Термины расширяемость и масштабируемость иногда используют как синонимы, но это неверно — каждый из них имеет чётко определённое самостоятельное значение. Расширяемость (extensibility) означает возможность сравнительно лёгкого добавления отдельных элементов сети (пользователей, компьютеров, приложений, служб), наращивания длины сегментов сети и замены существующей аппаратуры более мощной. При этом принципиально важно, что легкость расширения системы иногда может обеспечиваться в некоторых весьма ограниченных пределах. Например, локальная сеть Ethernet, построенная на основе одного сегмента толстого коаксиального кабеля, обладает хорошей расширяемостью, в том смысле, что позволяет легко подключать новые станции. Однако такая сеть имеет ограничение на число станций — оно не должно превышать 30–40. Хотя сеть допускает физическое подключение к сегменту и большего числа станций (до 100), но при этом чаще всего резко снижается производительность сети. Наличие такого ограничения и является признаком плохой масштабируемости системы при хорошей расширяемости. Масштабируемость (scalability) означает, что сеть позволяет наращивать количество узлов и протяжённость связей в очень широких пределах, при этом производительность сети не ухудшается. Для обеспечения масштабируемости сети приходится применять дополнительное коммуникационное оборудование и специальным образом структурировать сеть. Например, хорошей масштабируемостью обладает многосегментная сеть, построенная с использованием коммутаторов и маршрутизаторов и имеющая иерархическую структуру связей. Такая сеть может включать несколько тысяч компьютеров и при этом обеспечивать каждому пользователю сети нужное качество обслуживания.
\textit{(Источник: Олифер, гл. 2.)}

\paragraph{Прозрачность, управляемость и совместимость.}
Прозрачность (transparency) сети достигается в том случае, когда сеть представляется пользователям не как множество отдельных компьютеров, связанных между собой сложной системой кабелей, а как единая традиционная вычислительная машина с системой разделения времени. Известный лозунг компании Sun Microsystems: «Сеть — это компьютер» — говорит именно о такой прозрачной сети. Прозрачность может быть достигнута на двух различных уровнях — на уровне пользователя и на уровне программиста. На уровне пользователя прозрачность означает, что для работы с удалёнными ресурсами он использует те же команды и привычные ему процедуры, что и для работы с локальными ресурсами. На программном уровне прозрачность заключается в том, что приложению для доступа к удалённым ресурсам требуются те же вызовы, что и для доступа к локальным ресурсам. Управляемость сети подразумевает возможность централизованно контролировать состояние основных элементов сети, выявлять и разрешать проблемы, возникающие при работе сети, выполнять анализ производительности и планировать развитие сети. В идеале средства управления сетями представляют собой систему, осуществляющую наблюдение, контроль и управление каждым элементом сети — от простейших до самых сложных устройств, при этом такая система рассматривает сеть как единое целое, а не как разрозненный набор отдельных устройств. Совместимость, или интегрируемость, означает, что сеть способна включать в себя самое разнообразное программное и аппаратное обеспечение, то есть в ней могут сосуществовать различные операционные системы, поддерживающие разные стеки коммуникационных протоколов, а также аппаратные средства и приложения от разных производителей. Сеть, состоящая из разнотипных элементов, называется неоднородной, или гетерогенной, а если гетерогенная сеть работает без проблем, то она является интегрированной. Основной путь построения интегрированных сетей — использование модулей, выполненных в соответствии с открытыми стандартами и спецификациями.
\textit{(Источник: Олифер, гл. 2.)}

\paragraph{Качество обслуживания.}
В общем случае качество обслуживания (Quality of Service, QoS) определяет вероятностные оценки выполнения тех или иных требований, предъявляемых к сети приложениями или пользователями. Например, при передаче голосового трафика через сеть под качеством обслуживания чаще всего понимают гарантии того, что голосовые пакеты будут доставляться сетью с задержкой не более \(N\) мс, при этом вариация задержки не превысит \(M\) мс, и эти характеристики станут выдерживаться сетью с вероятностью 0,95 на определённом временном интервале. То есть приложению, которое передаёт голосовой трафик, важно, чтобы сеть гарантировала соблюдение именно этого набора характеристик качества обслуживания. Файловому сервису нужны гарантии средней полосы пропускания и расширения её на небольших интервалах времени до некоторого максимального уровня для быстрой передачи пульсаций. В идеале сеть должна гарантировать особые параметры качества обслуживания, сформулированные для каждого отдельного приложения. Однако по понятным причинам разрабатываемые и уже существующие механизмы QoS ограничиваются решением более простой задачи — гарантированием неких усреднённых требований, заданных для основных типов приложений. Чаще всего параметры, фигурирующие в разнообразных определениях качества обслуживания, регламентируют следующие показатели работы сети: скорость передачи данных; задержки передачи пакетов; уровень потерь и искажений пакетов. Качество обслуживания гарантируется для некоторого потока данных. Напомним, что поток данных — это последовательность пакетов, имеющих некоторые общие признаки, например адрес узла-источника, информация, идентифицирующая тип приложения (номер порта TCP/UDP) и т. п. К потокам применимы такие понятия, как агрегирование и дифференцирование. Так, поток данных от одного компьютера может быть представлен как совокупность потоков от разных приложений, а потоки от компьютеров одного предприятия агрегированы в один поток данных абонента некоторого поставщика услуг.
\textit{(Источник: Олифер, гл. 15.)}

\paragraph{Типы QoS и механизмы поддержки.}
Типы QoS различаются по степени «строгости», то есть по тому, насколько твёрдо сервис QoS может гарантировать обеспечение определённых значений характеристик QoS — пропускной способности, задержек, вариаций задержек, уровня потерь пакетов и т. п. Можно выделить три типа служб QoS. Сервис с максимальными усилиями, который также можно назвать отсутствием QoS, обеспечивает взаимодействие конечных узлов без каких бы то ни было гарантий. Типичные представители таких услуг — классические сети Ethernet или IP, которые не делают никаких различий между пакетами отдельных пользователей и приложений и обслуживают эти пакеты на основании принципа FIFO. Сервис с предпочтением (называемый также «мягким» сервисом QoS) — некоторые типы трафика обслуживаются лучше, чем остальные. Подразумеваются более быстрая обработка, в среднем больше пропускной способности и меньше потерь данных. Это статистическое предпочтение, а не численно выраженные гарантии. Гарантированный сервис (называемый также «жёстким», или «истинным», сервисом QoS) даёт статистические численные гарантии различным потокам трафика. Обычно такой вид QoS основан на предварительном резервировании сетевых ресурсов для каждого из потоков, получивших гарантии обслуживания. Механизмы поддержки качества обслуживания сами по себе не создают пропускной способности. Сеть не может дать больше того, что имеет. Механизмы QoS только управляют распределением имеющейся пропускной способности в соответствии с требованиями приложений и настройками сети. Для поддержки QoS применяются кондиционирование трафика, классификация, профилирование, формирование трафика, приоритетное и взвешенное обслуживание очередей, резервирование ресурсов и протоколы сигнализации. Реализация в компьютерных сетях механизмов поддержки QoS является сравнительно новой тенденцией. Долгое время компьютерные сети существовали без таких механизмов, и это объясняется в основном двумя причинами. Во-первых, большинство приложений, выполняемых в сети, были нетребовательными. Для таких приложений задержки пакетов или отклонения средней пропускной способности в достаточно широком диапазоне не приводили к значительной потере функциональности. Примерами нетребовательных приложений являются наиболее распространённые в сетях 80-х годов приложения электронной почты или удалённого копирования файлов. Во-вторых, сама пропускная способность 10-мегабитных сетей Ethernet во многих случаях не являлась дефицитом. В результате большинство сетей работало с тем качеством транспортного обслуживания, которое обеспечивало потребности приложений. Правда, никаких гарантий относительно нахождения задержек пакетов или пропускной способности, с которой пакеты передаются между узлами, в определённых пределах эти сети не давали. Более того, при временных перегрузках сети, когда значительная часть компьютеров одновременно начинала передавать данные с максимальной скоростью, задержки и пропускная способность становились такими, что работа приложений давала сбой — шла слишком медленно, с прерываниями сеансов и т. п. Транспортный сервис, который предоставляли такие сети, получил название best effort, то есть сервис «с максимальными усилиями». Сеть старается обработать поступающий трафик как можно быстрее, но при этом никаких гарантий относительно результата своих усилий не даёт.
\textit{(Источник: Олифер, гл. 15.)}

\subsubsection{Краткое резюме подраздела}
Информационно-вычислительная сеть является распределённой вычислительной системой, в которой группа компьютеров согласованно решает взаимосвязанные задачи, обмениваясь данными в автоматическом режиме. К распределённым системам также относятся мультипроцессоры и кластеры. В сетях связи между компьютерами слабее, чем в мультипроцессорах и кластерах, а автономность узлов выше. Разделение ресурсов в сети строится на клиент-серверной модели и реализуется через сетевые службы. Распределённые приложения состоят из нескольких частей, выполняющихся на разных компьютерах, и в полной мере используют возможности распределённой обработки. Для повышения надёжности и доступности применяются репликация данных и удалённые вызовы процедур. Важными свойствами сети являются прозрачность, управляемость, совместимость, расширяемость, масштабируемость, надёжность, безопасность и качество обслуживания. Механизмы QoS позволяют управлять распределением пропускной способности и обеспечивать различные уровни гарантий для разных типов трафика. Совместное применение этих принципов и механизмов обеспечивает построение эффективных, надёжных и масштабируемых информационно-вычислительных сетей.
\textit{(Обобщение по: Олифер, гл. 1–2, 15.)}

\end{document}